% =========================================================
% Template TCC - UNISALES
% Elaborado por: Gustavo Cardoso Rodrigues
% Tema: Gestão de Chamados Internos entre Áreas
% =========================================================

\documentclass[a4paper,12pt]{article}

\input{preambulo}

\begin{document}

% ================= LOGOS =================
\begin{center}
\includegraphics[width=0.4\textwidth]{Imagens/ciencia_pratica.png}
\hspace{1cm}
\includegraphics[width=0.25\textwidth]{Imagens/Unisales.png}
\end{center}

\vspace{1cm}

% ================= TÍTULO E AUTORES =================
\begin{center}

\textbf{Low code na gestão de chamados internos entre áreas (Provisório)}

\vspace{0.3cm}

\textit{LOW CODE IN INTERNAL TICKET MANAGEMENT BETWEEN DEPARTMENTS (PROVISIONAL)}

\vspace{0.8cm}

Gustavo Cardoso Rodrigues\footnote{Centro Universitário Salesiano (UNISALES). Vitória/ES, Brasil.}

\end{center}

\vspace{1cm}

% ================= RESUMO =================
\textbf{RESUMO:} A ausência de um sistema formalizado para o
gerenciamento de chamados internos entre áreas gera perda de
informações, impossibilidade de rastreabilidade e ausência de
indicadores de desempenho nas organizações que ainda conduzem
esse fluxo por canais informais. Este trabalho tem como
objetivo analisar o processo atual de gestão de chamados
internos entre áreas, levantar requisitos, redesenhar e
desenvolver um sistema de gestão de chamados, utilizando
plataformas do ecossistema Microsoft 365, em uma empresa de
grande porte do setor de serviços de transporte, localizada
no Estado do Espírito Santo. A solução é estruturada em torno
de um portal no SharePoint como ponto único de entrada, um
aplicativo em Power Apps para abertura padronizada dos
chamados, Power Automate para automação do fluxo e das
notificações, Microsoft List como repositório centralizado e
Power BI para geração de indicadores gerenciais. A pesquisa é
de natureza aplicada, com abordagem qualitativa, conduzida
como estudo de caso único na organização em que o autor atua
como assistente de dados, tendo a Engenharia de Processos e o
\textit{Business Process Management} (BPM) como referenciais
para a análise do processo atual (AS-IS) e para o redesenho
do processo futuro (TO-BE) que precede o desenvolvimento da
solução. Os resultados esperados incluem eliminação da
comunicação informal entre setores, rastreabilidade integral
dos chamados, controle de SLA por prioridade e geração de
indicadores de volume, tempo médio de resolução e causa raiz
dos problemas.
 
\textbf{Palavras-chave:} Gestão de Chamados Internos;
Engenharia de Processos; Automação de Processos; Low-Code;
Microsoft 365; Business Intelligence; BPM; SLA.
 
\vspace{0.08cm}
 
\textbf{ABSTRACT:} The absence of a formal system for managing
internal tickets between departments causes information loss,
lack of traceability, and absence of performance indicators in
organizations that still handle this flow through informal
channels. This study aims to analyze the current process of
managing internal tickets between departments, gather
requirements, redesign the process, and develop a ticket
management system using Microsoft 365 ecosystem platforms, in
a large company in the transportation services sector, located
in the state of Espírito Santo, Brazil. The solution is built
around a SharePoint portal as the single point of entry, a
Power Apps application for standardized ticket opening, Power
Automate for workflow and notification automation, Microsoft
List as a centralized repository, and Power BI for managerial
indicators. The research is applied in nature, with a
qualitative approach, conducted as a single case study in the
organization where the author works as a data assistant, using
Process Engineering and Business Process Management (BPM) as
references for analyzing the current process (AS-IS) and
redesigning the future process (TO-BE) that precedes the
development of the solution. Expected results include
elimination of informal inter-departmental communication, full
ticket traceability, SLA control by priority, and generation of
indicators for volume, average resolution time, and root cause
analysis.
 
\textbf{Keywords:} Internal Ticket Management; Process
Engineering; Process Automation; Low-Code; Microsoft 365;
Business Intelligence; BPM; SLA.
 
\vspace{1cm}

% ================= CONTEÚDO =================
\section{Introdução}

A transformação digital tem exigido das organizações não
apenas a adoção de novas ferramentas, mas a revisão de seus
processos internos, especialmente aqueles que ainda dependem
de comunicação informal entre setores. Quando problemas
recorrentes são reportados por aplicativos de mensagens ou
e-mails não estruturados, sem protocolo e sem registro
centralizado, a organização perde rastreabilidade, dificulta
a gestão e impossibilita a geração de indicadores de
desempenho \cite{Ramos2024}. A literatura recente sobre
transformação digital é enfática ao afirmar que a
digitalização de um fluxo manual, quando realizada sem a
revisão estrutural do processo que lhe deu origem, tende
apenas a reproduzir em ambiente eletrônico as mesmas
ineficiências já existentes \cite{SantosCaldeiraPacheco2026}.
Este trabalho parte dessa premissa: antes de desenvolver uma
aplicação, é necessário compreender e redesenhar o processo
que ela pretende sustentar.

Assim, a lógica de construção desta pesquisa segue uma
sequência específica --- do problema organizacional à
automação --- que também organiza a estrutura dos capítulos
deste trabalho: identificação do problema organizacional,
análise e entendimento do processo existente, identificação
dos gargalos e necessidades, levantamento de requisitos,
redesenho do processo e, somente então, desenvolvimento da
solução tecnológica e automação. Essa sequência é sustentada
pela integração entre a Engenharia de Processos, o
\textit{Business Process Management} (BPM), a Engenharia de
Software, a Engenharia de Requisitos e a abordagem
\textit{low-code}, discutidas na Base Teórica deste trabalho.

Em organizações que ainda não dispõem de um canal único de
abertura de chamados internos, problemas que exigem análise e
resolução por diferentes áreas são comunicados de forma não
padronizada. Não há protocolo gerado automaticamente, não há
controle de prazo e não há visibilidade gerencial sobre o
volume, a prioridade ou o tempo médio de resolução dessas
demandas. A ausência de um canal formalizado de entrada expõe
a organização ao retrabalho, à perda de informações e à
impossibilidade de identificar causas raiz dos problemas
recorrentes \cite{Santos2023}.

Diante desse cenário, este trabalho realiza a análise do
processo atual (AS-IS), o levantamento de requisitos, a
modelagem do processo futuro (TO-BE), o desenvolvimento e a
implantação de um fluxo de gestão de chamados internos entre
áreas, tendo como ponto de entrada um site desenvolvido no
\textbf{SharePoint}, responsável por centralizar a abertura
das solicitações. Os chamados registrados nesse site
constituem a entrada de um processo automatizado desenvolvido
com base no ecossistema Microsoft~365, contemplando o registro
estruturado da solicitação, a padronização do fluxo entre as
áreas envolvidas e a geração de indicadores gerenciais de
acompanhamento \cite{Lopes2024}.

O objetivo geral deste trabalho é analisar o processo atual
de comunicação de chamados internos entre áreas, levantar os
requisitos, redesenhar e modelar o processo, e desenvolver e
implantar um sistema de gestão de chamados internos, tendo
como ponto de entrada um site no SharePoint e utilizando
ferramentas do ecossistema Microsoft 365, de modo a
padronizar a abertura das solicitações, eliminar a
comunicação informal e viabilizar a geração de indicadores
gerenciais de desempenho do processo.

% ─── JUSTIFICATIVA ────────────────────────────────────────
\subsection{Justificativa do Tema}

A escolha do tema fundamenta-se em quatro dimensões
complementares: relevância prática imediata, aderência à
disciplina de Engenharia de Processos e de Gestão de
Processos de Negócio, aderência às competências da
Engenharia de Software e da Engenharia de Requisitos, e
contribuição ao conhecimento acadêmico aplicado.

Do ponto de vista prático, a organização estudada não dispõe
de qualquer mecanismo formalizado para abertura,
acompanhamento e encerramento de chamados internos entre
áreas. O autor possui conhecimento direto das dificuldades
operacionais do processo e acesso às informações necessárias
ao seu diagnóstico e ao levantamento de requisitos
\cite{Gil2002}. A ausência de um canal formalizado de abertura
de chamados gera custos ocultos decorrentes de retrabalho,
perda de informações em canais informais e impossibilidade de
gestão baseada em dados.

Do ponto de vista da Engenharia de Processos e do BPM, o tema
exige que a solução tecnológica seja precedida pela
compreensão do processo existente e pelo seu redesenho
estruturado, superando a lógica de que a automação, por si
só, resolve problemas de gestão \cite{SantosCaldeiraPacheco2026}.
Do ponto de vista da Engenharia de Software e da Engenharia
de Requisitos, o tema mobiliza competências nucleares do
curso: especificação de requisitos a partir da modelagem do
processo de negócio \cite{CarvalhoJunior2022}, arquitetura de
solução orientada a serviços e avaliação de qualidade com
usuários. A construção de um sistema de chamados com controle
de status, prioridade, SLA, causa raiz e rastreabilidade,
integrado a um painel gerencial, constitui um projeto de
engenharia completo, com requisitos funcionais e não
funcionais bem definidos, mesmo em ambiente \textit{low-code}
\cite{Alves2022}.

Do ponto de vista acadêmico, a pesquisa contribui com
evidência empírica sobre a aplicação integrada da Engenharia
de Processos, do BPM, da Engenharia de Requisitos e do
\textit{low-code} na modelagem de um processo de gestão de
chamados internos, a partir de um site no SharePoint como
ponto único de entrada, integrado ao ecossistema Microsoft
365 --- recorte pouco explorado de forma conjunta na
literatura nacional recente, constituindo contribuição
original ao campo \cite{Santos2023}.

% ─── PROBLEMA ─────────────────────────────────────────────
\subsection{Problema da Pesquisa}

Os chamados internos entre áreas da organização são
comunicados por aplicativos de mensagens e e-mails sem
padronização, sem geração de protocolo, sem controle de prazo
e sem registro centralizado. A área responsável pela análise
recebe a demanda sem garantia de completude das informações.
A gestão não possui indicadores sobre volume, prioridade,
tempo de resolução ou causa raiz dos problemas. Chamados são
esquecidos ou duplicados sem que haja mecanismo de alerta ou
escalonamento.

Diante desse contexto, formula-se a seguinte questão de
pesquisa:

\textit{Como estruturar, a partir da análise do processo
atual e de um site no SharePoint como ponto único de entrada,
o levantamento de requisitos, o redesenho, a modelagem, o
desenvolvimento e a implantação de um processo automatizado
de gestão de chamados internos entre áreas, capaz de
padronizar a abertura, garantir rastreabilidade, controlar
SLA por prioridade, registrar causa raiz e gerar indicadores
gerenciais de volume, tempo médio de resolução e percentual
de chamados concluídos dentro do prazo?}

% ─── HIPÓTESE ─────────────────────────────────────────────
\subsection{Hipótese}

A análise do processo atual, o levantamento de requisitos, o
redesenho segundo os princípios do BPM, o desenvolvimento e a
implantação de um processo de gestão de chamados internos
entre áreas --- tendo como ponto de entrada um site no
SharePoint e como componentes de automação a Microsoft List,
o Power Automate e o Power BI --- permitem obter uma solução
capaz de eliminar a comunicação informal entre setores,
garantir rastreabilidade integral dos chamados, controlar SLA
por prioridade, registrar causa raiz e ação corretiva, e
gerar indicadores gerenciais de volume, tempo médio de
resolução e percentual de atendimentos dentro do prazo,
contribuindo para a melhoria mensurável da eficiência
operacional das áreas envolvidas.

% ─── ESTADO DA ARTE ───────────────────────────────────────
\subsection{Estado da Arte}

O levantamento realizado para este trabalho abrangeu
produções em língua portuguesa, publicadas nos últimos cinco
anos e indexadas no Google Scholar, SciELO e Portal de
Periódicos CAPES, com os descritores: \textit{engenharia de
processos}, \textit{BPM}, \textit{engenharia de requisitos},
\textit{elicitação de requisitos}, \textit{low-code},
\textit{Microsoft List}, \textit{Power Automate},
\textit{automação de processos}, \textit{SLA} e
\textit{indicadores de desempenho}.

No campo da Engenharia de Processos e da transformação
digital, \cite{SantosCaldeiraPacheco2026} demonstram que a
transformação digital efetiva depende da revisão e do
redesenho estruturado dos processos organizacionais, e não da
simples informatização de fluxos preexistentes, resultado que
orienta diretamente a decisão deste trabalho de tratar o
redesenho do processo como etapa anterior ao desenvolvimento.
No campo da gestão de processos, \cite{Ramos2024} demonstram
que a disponibilização de dados gerenciais estruturados é o
resultado de maior valor percebido após iniciativas de BPM em
instituições de médio e grande porte. Em linha com esse
resultado, \cite{Lopes2024} evidenciam que a digitalização de
processos com BPM e ferramentas do ecossistema Microsoft
resultou em repositório único de dados com rastreabilidade
completa, eliminando a comunicação informal entre setores.

No campo da Engenharia de Requisitos, \cite{CarvalhoJunior2022}
demonstram que a derivação de requisitos de software a partir
da modelagem prévia do processo de negócio produz requisitos
mais alinhados às metas organizacionais do que técnicas
tradicionais de elicitação isoladas, resultado que reforça a
opção metodológica deste trabalho de conduzir o levantamento
de requisitos junto a perfis que representam diferentes níveis
de atuação no processo.

No campo das plataformas \textit{low-code},
\cite{Alves2022} identificam redução de até 60\% no tempo de
desenvolvimento de soluções corporativas com o paradigma
\textit{low-code} em comparação ao desenvolvimento
tradicional, destacando o ecossistema Microsoft 365 como
plataforma de maior penetração no mercado nacional.
\cite{Santos2023} concluem que a combinação de BPM com
\textit{low-code} é especialmente eficaz na construção de
repositórios e painéis de acompanhamento, pois permite que as
regras de negócio sejam configuradas pelas próprias áreas
proprietárias dos processos. \cite{Costa2022} demonstram, em
estudo de caso aplicado, que a adoção de ferramentas
\textit{low-code} resultou em rastreabilidade completa e
redução de 45\% no tempo de processamento de solicitações
internas.

No campo da automação de processos, \cite{ReisFaria2024}
evidenciam, em revisão sistemática sobre automação robótica
de processos, que os ganhos de padronização e redução de
retrabalho são mais consistentes quando a automação é
aplicada sobre um processo já revisado. No campo da gestão de
SLA, \cite{Malheiros2024} concluem que o SLA é fundamental
para estabelecer expectativas claras entre áreas, promovendo
qualidade, transparência e rastreabilidade na resolução de
chamados, e destacam que a automação do controle de prazos é
condição para que o SLA deixe de ser apenas um documento
formal.

O levantamento evidencia uma lacuna na literatura nacional
quanto à aplicação integrada da Engenharia de Processos, do
BPM e da Engenharia de Requisitos na construção de processos
de gestão de chamados internos entre áreas a partir do
ecossistema Microsoft 365, com um site no SharePoint como
ponto único de entrada. O presente trabalho busca preencher
essa lacuna com a análise do processo atual, o levantamento
de requisitos, o redesenho, o desenvolvimento e a implantação
do processo, documentados de forma estruturada.

% ─── OBJETIVOS ────────────────────────────────────────────
\subsection{Objetivos}

\subsubsection{Objetivo Geral}

Analisar o processo atual de gestão de chamados internos
entre áreas, levantar os requisitos, redesenhar e modelar o
processo, e desenvolver e implantar um sistema de gestão de
chamados internos, tendo como ponto de entrada um site no
SharePoint e utilizando plataformas do ecossistema Microsoft
365, de modo a padronizar a abertura, garantir
rastreabilidade, controlar SLA por prioridade, registrar
causa raiz e gerar indicadores gerenciais de desempenho do
processo.

\subsubsection{Objetivos Específicos}

\begin{itemize}
    \item Analisar o processo atual (AS-IS) de comunicação de
    chamados internos entre áreas, identificando participantes,
    atividades manuais e pontos de falha;

    \item Levantar e especificar os requisitos funcionais e
    não funcionais do processo por meio de um formulário
    estruturado, aplicado aos profissionais responsáveis pelo
    processo em diferentes níveis de atuação (Analista,
    Supervisora e Gerente), com perguntas direcionadas
    conforme o nível de responsabilidade de cada participante;

    \item Redesenhar o processo (TO-BE) a partir dos
    requisitos levantados, definindo os campos obrigatórios de
    abertura no site do SharePoint, os status possíveis
    (Aberto, Em análise, Concluído), os critérios de
    prioridade (Alta, Média, Baixa) e o tempo de SLA associado
    a cada nível de prioridade;

    \item Especificar o repositório centralizado a ser
    integrado ao site do SharePoint, com os campos: Protocolo
    N\textsuperscript{o}, Tipo de chamado (Massivo ou Interno),
    Assunto, Problema, Quantidade de protocolos relacionados,
    Área Responsável, Prioridade, Evidências, Solicitante,
    Status, Causa raiz, Ação corretiva, SLA (contagem
    automática) e Data da solução;

    \item Desenvolver o fluxo automatizado responsável por
    registrar cada chamado a partir do site do SharePoint,
    gerar protocolo único, enviar notificação à área
    responsável com link de acesso direto ao chamado e
    realizar a contagem automática do SLA;

    \item Desenvolver e implantar o painel gerencial com os
    indicadores: chamados abertos, em análise e concluídos por
    período; tempo médio de resolução; percentual de chamados
    dentro do SLA; distribuição por área responsável e por
    prioridade; e base histórica com rastreabilidade completa
    por protocolo;

    \item Validar o novo modelo do processo junto aos
    participantes envolvidos, a partir do fluxo implantado e
    do painel gerencial em uso.
\end{itemize}

% ─── BASE TEÓRICA ─────────────────────────────────────────
\section{Base Teórica}

Esta seção apresenta a fundamentação teórica que sustenta o
trabalho, organizada de modo a evidenciar a lógica de
construção da pesquisa: o problema organizacional é
compreendido a partir da Engenharia de Processos, o
redesenho do fluxo é conduzido segundo a disciplina do
\textit{Business Process Management} (BPM), a transformação
desse redesenho em solução computacional é orientada pela
Engenharia de Software e pela Engenharia de Requisitos, a
implementação é viabilizada pela abordagem \textit{low-code},
e o resultado operacional é analisado à luz da automação de
processos. Dessa forma, a tecnologia não é apresentada como
ponto de partida do trabalho, mas como consequência de um
processo de análise e redesenho previamente conduzido.

\subsection{Engenharia de Processos e Transformação
Organizacional}

A Engenharia de Processos constitui a etapa inicial da
transformação digital, na medida em que estabelece que a
implantação de qualquer solução tecnológica deve decorrer de
uma necessidade operacional previamente identificada, e não
o contrário. Sob essa perspectiva, compreender o processo
existente --- seus participantes, atividades manuais, pontos
de falha e oportunidades de melhoria --- é condição prévia
para que a automação produza ganhos reais de padronização,
rastreabilidade e governança \cite{SantosCaldeiraPacheco2026}.

Um dos instrumentos centrais dessa abordagem é a
representação do processo em dois estados complementares.
O processo atual, ou \textbf{AS-IS}, descreve o fluxo antes
da intervenção, evidenciando comunicação descentralizada,
ausência de padronização, dependência de controles manuais e
baixa rastreabilidade das informações. O processo futuro, ou
\textbf{TO-BE}, representa o fluxo redesenhado após a
implantação da solução, contemplando abertura estruturada,
registro centralizado, definição clara de responsáveis,
automação de notificações e acompanhamento do ciclo de vida
de cada ocorrência \cite{SantosCaldeiraPacheco2026}. A
transição entre esses dois estados não decorre da simples
transposição do fluxo manual para um ambiente digital: quando
a digitalização ocorre sem revisão estrutural do processo, ela
tende a reproduzir, em formato eletrônico, as mesmas
ineficiências do modelo anterior
\cite{SantosCaldeiraPacheco2026}. O foco desta pesquisa,
portanto, não é demonstrar que um novo sistema foi criado,
mas evidenciar que o processo de gestão de chamados internos
foi analisado e redesenhado antes de sua automação.

\subsection{Gestão de Processos de Negócio (BPM)}

O \textit{Business Process Management} (BPM) fornece o
arcabouço metodológico que estrutura o ciclo de
identificação, modelagem, análise, redesenho, automação e
monitoramento de um processo organizacional
\cite{Ramos2024}. Trata-se de uma disciplina que orienta a
organização a superar a lógica funcional tradicional,
substituindo-a por uma visão orientada a processos, capaz de
atravessar diferentes áreas e gerar valor de ponta a ponta
\cite{SantosCaldeiraPacheco2026}.

No contexto de chamados internos entre áreas, o BPM orienta a
definição dos estados possíveis do chamado, as
responsabilidades de cada ator no fluxo e os critérios de
encerramento. \cite{Lopes2024} demonstraram que a
digitalização de processos com BPM e ferramentas do
ecossistema Microsoft resultou em repositório único de dados
com rastreabilidade completa, eliminando a comunicação
informal entre setores. Em sentido semelhante,
\cite{Ramos2024} evidenciam que a disponibilização de dados
gerenciais estruturados constitui o resultado de maior valor
percebido após iniciativas de BPM em instituições de médio e
grande porte, ao passo que
\cite{SantosCaldeiraPacheco2026} destacam que a efetividade
da transformação digital depende da revisão estruturada dos
processos e do fortalecimento das capacidades
organizacionais, e não apenas da incorporação de novas
ferramentas.

O projeto desenvolvido neste trabalho é relacionado às
etapas do ciclo de BPM conforme apresentado no
Quadro~\ref{quadro:bpm-projeto}.

\begin{table}[H]
\caption{Relação entre as etapas do BPM e o projeto desenvolvido}
\label{quadro:bpm-projeto}
\centering
\begin{tabular}{p{2.6cm} p{9.2cm}}
\hline
\textbf{Etapa do BPM} & \textbf{Aplicação no projeto} \\
\hline
Identificação & Análise do processo de comunicação de
chamados internos entre áreas. \\
\hline
Modelagem & Construção da representação do fluxo atual
(AS-IS) e do fluxo futuro (TO-BE). \\
\hline
Análise & Identificação dos gargalos operacionais a partir
do levantamento de requisitos junto aos participantes. \\
\hline
Redesenho & Definição do novo fluxo de atendimento, com
estados, responsáveis e regras de negócio. \\
\hline
Automação & Desenvolvimento dos componentes no ecossistema
Microsoft 365 (Power Apps, Microsoft List e Power
Automate). \\
\hline
Monitoramento & Construção do painel gerencial no Power BI
como base para indicadores de desempenho. \\
\hline
\end{tabular}
\raggedright
\par
{\fontsize{10}{12}\selectfont Fonte: Elaborado pelo autor,
adaptado de \cite{Ramos2024} e \cite{SantosCaldeiraPacheco2026}}
\end{table}

A utilização do BPM neste trabalho reforça que a tecnologia
foi aplicada como consequência da melhoria do processo, e
não como ponto inicial da solução.

\subsection{Engenharia de Software em Ambiente Low-Code}

A Engenharia de Software é a disciplina responsável por
transformar necessidades organizacionais em uma solução
tecnológica estruturada, por meio da aplicação sistemática de
métodos de análise, projeto, implementação e validação
\cite{Alves2022}. A adoção de uma plataforma \textit{low-code}
não elimina essa disciplina: ela modifica a forma de
implementação, substituindo a codificação manual extensiva
por configuração visual e componentes pré-construídos, mas
mantém a necessidade de análise de requisitos, definição de
arquitetura, estruturação de dados e validação do
comportamento esperado \cite{Mendes2024}.

Este trabalho relaciona explicitamente as etapas tradicionais
da Engenharia de Software às atividades desenvolvidas em
ambiente \textit{low-code}: o levantamento da necessidade
corresponde à pesquisa de campo conduzida junto aos três
perfis profissionais envolvidos no processo; a especificação
de requisitos corresponde à consolidação dos requisitos
funcionais e não funcionais; a definição de arquitetura
corresponde à escolha dos componentes do ecossistema
Microsoft 365 e à forma como eles se integram; a
implementação corresponde à construção do aplicativo, da
lista e dos fluxos de automação; e os testes e a validação
correspondem à verificação do funcionamento integrado e à
avaliação da solução junto aos participantes. Essa
correspondência evidencia que, mesmo em ambiente
\textit{low-code}, o projeto preserva as características de
um trabalho de engenharia, com requisitos funcionais e não
funcionais bem definidos.

\subsection{Engenharia de Requisitos}

A Engenharia de Requisitos é a área da Engenharia de Software
responsável por identificar, documentar, validar e gerenciar
as necessidades dos \textit{stakeholders} ao longo do ciclo
de desenvolvimento de um sistema. A qualidade dos requisitos
influencia diretamente a qualidade da solução produzida, pois
falhas introduzidas nessa etapa tendem a gerar retrabalho e
aumento de custos \cite{Alves2022}. Segundo
\cite{CarvalhoJunior2022}, a elicitação constitui uma das
fases mais críticas da Engenharia de Requisitos, pois a
qualidade de um software depende da correta compreensão das
necessidades dos usuários e da capacidade de alinhar os
requisitos gerados às metas do negócio --- alinhamento que,
segundo os mesmos autores, é favorecido quando os requisitos
são derivados a partir da modelagem prévia do processo de
negócio, aproximando a Engenharia de Requisitos da Engenharia
de Processos discutida anteriormente.

Neste trabalho, cada necessidade identificada junto aos
participantes da pesquisa foi conduzida por um percurso
específico até se tornar parte da solução, conforme
ilustrado no exemplo do Quadro~\ref{quadro:necessidade-requisito}.

\begin{table}[H]
\caption{Do problema organizacional ao requisito
implementado}
\label{quadro:necessidade-requisito}
\centering
\begin{tabular}{p{3.4cm} p{3.9cm} p{4.5cm}}
\hline
\textbf{Necessidade identificada} & \textbf{Requisito
funcional} & \textbf{Implementação} \\
\hline
Dificuldade para acompanhar quem está responsável pelo
chamado. & O sistema deve permitir identificar o responsável
pela tratativa. & Campo de responsável pela tratativa
integrado ao fluxo de atendimento na Microsoft List. \\
\hline
\end{tabular}
\raggedright
\par
{\fontsize{10}{12}\selectfont Fonte: Elaborado pelo autor}
\end{table}

Os requisitos de um sistema dividem-se, de forma clássica, em
requisitos funcionais, que descrevem os serviços que o
sistema deve oferecer, e requisitos não funcionais, que
estabelecem restrições de qualidade como segurança,
rastreabilidade, usabilidade e disponibilidade
\cite{Alves2022}. No contexto deste trabalho, os requisitos
funcionais contemplam a abertura padronizada de chamados, a
geração automática de protocolo, o controle de status e a
geração de indicadores gerenciais; os requisitos não
funcionais incluem a operação dentro do ecossistema Microsoft
365 já licenciado pela organização e o acesso diferenciado por
perfil de usuário \cite{Lopes2024}.

Entre as técnicas de elicitação de requisitos, o uso de
questionários e formulários estruturados destaca-se por
permitir a coleta padronizada de informações junto a
diferentes perfis de \textit{stakeholders}, com menor custo
operacional em comparação a técnicas presenciais
\cite{Gil2002}. A segmentação das perguntas conforme o papel
do respondente no processo --- nível operacional, de
supervisão ou de gestão --- favorece a triangulação das
informações coletadas, permitindo validar o mesmo fluxo sob
diferentes perspectivas organizacionais.

\subsection{Plataformas Low-Code como Estratégia Tecnológica}

As plataformas \textit{low-code} constituem uma estratégia
tecnológica para implementação de soluções corporativas,
especialmente em organizações que não dispõem de equipes de
desenvolvimento dedicadas. O paradigma \textit{low-code}
minimiza a necessidade de codificação manual por meio de
interfaces visuais e componentes pré-configurados, permitindo
que profissionais com menor especialização técnica construam
e mantenham aplicações funcionais \cite{Alves2022}.
\cite{Santos2023} destacam que essa característica tem
viabilizado o fenômeno dos \textit{citizen developers} ---
colaboradores das áreas de negócio que assumem papel ativo na
construção de soluções digitais sem depender exclusivamente
de equipes de TI --- e que a combinação entre BPM e
\textit{low-code} é especialmente eficaz na construção de
repositórios e painéis de acompanhamento, pois permite que as
regras de negócio sejam configuradas pelas próprias áreas
proprietárias dos processos. \cite{Costa2022} demonstraram,
em estudo de caso aplicado, que a adoção de ferramentas
\textit{low-code} resultou em rastreabilidade completa e
redução de 45\% no tempo de processamento de solicitações
internas, com desenvolvimento concluído em semanas.

Este trabalho não pretende apresentar o \textit{low-code}
como substituto da programação tradicional, pois não há, na
literatura consultada, métricas suficientes para sustentar
tal comparação. O foco adotado é demonstrar como a plataforma
possibilitou implementar uma solução integrada ao ambiente
corporativo já existente, reduzindo barreiras técnicas e
permitindo a criação de uma aplicação alinhada às
necessidades do processo redesenhado. A arquitetura adotada é
composta por quatro componentes do ecossistema Microsoft 365.
O primeiro é o \textbf{Power Apps}, interface de abertura e
acompanhamento do chamado. O segundo é a \textbf{Microsoft
List}, evolução das listas do SharePoint Online, que funciona
como repositório centralizado das informações, herdando
recursos de controle de versão, permissões por perfil e
integração nativa com o Power Automate e o Power BI
\cite{Mendes2024}. O terceiro é o \textbf{Power Automate},
responsável pela execução das regras de negócio e das
notificações automáticas. O quarto é o \textbf{Power BI},
responsável pela análise gerencial e geração de indicadores.
Essa arquitetura reflete o modelo de integração nativa entre
componentes do Microsoft 365, sem necessidade de conectores
externos ou infraestrutura adicional \cite{Costa2022}.

\subsection{Automação de Processos}

A automação de processos é o mecanismo pelo qual regras de
negócio previamente definidas passam a ser executadas de
forma padronizada, reduzindo a dependência de atividades
manuais para comunicação, direcionamento e acompanhamento de
demandas \cite{ReisFaria2024}. Estudos sobre automação
robótica de processos (RPA) em contextos organizacionais
distintos do estudado neste trabalho indicam ganhos
recorrentes de padronização e redução de retrabalho quando a
automação é aplicada sobre um processo previamente revisado,
e não sobre um fluxo manual simplesmente transposto para o
ambiente digital \cite{ReisFaria2024}.

No processo analisado, a automação substitui uma sequência de
etapas dependentes de ação humana --- identificação da
necessidade, envio de comunicação e acompanhamento manual do
retorno --- por uma sequência orientada a eventos, na qual a
alteração de um registro dispara automaticamente a regra
correspondente, a notificação ao responsável e a atualização
do histórico. Essa transformação está alinhada ao que
\cite{Lopes2024} descrevem como resultado da digitalização de
processos com ferramentas do ecossistema Microsoft: repositório
único de dados, com rastreabilidade completa e eliminação da
comunicação informal entre setores. A automação, portanto, não
é tratada neste trabalho como mecanismo isolado de registro,
mas como consequência do redesenho do processo discutido nas
seções anteriores.

\subsection{SLA na Gestão de Chamados Internos}

O Acordo de Nível de Serviço (\textit{Service Level
Agreement} --- SLA) é o instrumento que formaliza prazos,
critérios de qualidade e responsabilidades no atendimento de
demandas internas. \cite{Malheiros2024} concluem que o SLA é
fundamental para estabelecer expectativas claras entre
provedores e usuários de serviços, promovendo qualidade,
transparência e resolução eficaz de disputas, e destacam que
a automação do controle de prazos --- com geração de alertas
quando um chamado se aproxima do vencimento --- é condição
para que o SLA deixe de ser apenas um documento formal e passe
a ser um instrumento ativo de gestão. No contexto deste
trabalho, o SLA é incorporado como requisito não funcional e
como regra de negócio associada à prioridade de cada chamado,
integrando-se ao ciclo de monitoramento previsto pelo BPM.

\subsection{Comunicação Interna e Gestão de Informações entre
Áreas}

A comunicação interna é reconhecida na literatura como fator
crítico para a eficiência operacional e para o desempenho
organizacional. \cite{Souza2023} concluem, em revisão de
literatura, que grande parte das organizações tem dificuldade
em manter uma comunicação interna eficiente, optando, na
maioria das vezes, por canais que não garantem registro,
completude ou rastreabilidade das informações trocadas entre
setores. A substituição de canais informais por sistemas
estruturados, com registro automático e notificação
padronizada, reduz erros, elimina retrabalho e melhora o
desempenho das equipes envolvidas \cite{Lopes2024}.

No contexto específico da gestão de chamados internos, a
formalização do fluxo por meio de um canal estruturado
representa a passagem de um modelo reativo --- no qual o
problema é comunicado de forma \textit{ad hoc} --- para um
modelo proativo, no qual cada chamado gera automaticamente um
protocolo, uma notificação à área responsável e um registro
rastreável desde a abertura até o encerramento. Essa
transição é o elo final entre os seis eixos teóricos
apresentados nesta seção: o problema organizacional
compreendido pela Engenharia de Processos é redesenhado sob a
disciplina do BPM, especificado pela Engenharia de Requisitos,
implementado pela Engenharia de Software em ambiente
\textit{low-code} e operacionalizado pela automação, com a
comunicação estruturada entre áreas como resultado observável
da transformação.

\section{Metodologia}

\subsection{Caracterização da Pesquisa}

Quanto à natureza, a presente pesquisa é classificada como
\textbf{aplicada}, pois busca gerar conhecimento voltado à
solução de um problema real identificado em uma organização
\cite{Gil2002}: o objetivo é analisar o processo atual,
levantar os requisitos, redesenhar e modelar o processo, e
propor a automação de um sistema de gestão de chamados
internos entre áreas, tendo como ponto de entrada um site no
SharePoint e como componentes de automação ferramentas do
ecossistema Microsoft 365.

Quanto à abordagem, a pesquisa é \textbf{qualitativa}, uma
vez que busca compreender as necessidades do contexto
organizacional a partir da perspectiva dos profissionais
diretamente envolvidos no processo, transformando essas
necessidades em requisitos para o redesenho e o
desenvolvimento de uma solução computacional \cite{Gil2002}.

Quanto ao método, a pesquisa caracteriza-se como
\textbf{estudo de caso}, pois analisa um processo específico
--- a gestão de chamados internos entre áreas --- dentro de
uma organização determinada, e desenvolve uma solução
direcionada ao contexto estudado \cite{Gil2002}. Quanto aos
objetivos, a pesquisa é \textbf{exploratória e descritiva}:
exploratória, por investigar a aplicação integrada da
Engenharia de Processos, do BPM e do ecossistema Microsoft
365 na gestão de chamados internos entre áreas; descritiva,
por documentar os requisitos levantados junto aos
participantes, a arquitetura desenvolvida e os resultados
obtidos na avaliação da solução \cite{Santos2023}.

O delineamento adotado é o de \textbf{pesquisa de campo}, uma
vez que a coleta de dados ocorre diretamente no ambiente
organizacional, por meio de um formulário estruturado (Google
Forms) aplicado aos profissionais responsáveis pelo processo
estudado \cite{Gil2002}. Essa escolha permite capturar as
necessidades reais do processo a partir de quem o opera
diariamente, garantindo que os requisitos da solução reflitam
a realidade da organização, além de possibilitar a validação
de todo o fluxo interno da futura automação a partir de
diferentes níveis de atuação dos respondentes.

\subsection{Participantes da Pesquisa}

Os participantes foram selecionados por critério de
intencionalidade, considerando o papel de cada um no processo
de gestão de chamados internos da organização. São três
profissionais da área de Qualidade:

\begin{itemize}
    \item \textbf{Gerente da Área} --- responsável por
    validar a aderência do fluxo proposto aos objetivos
    estratégicos, os benefícios esperados, a viabilidade da
    solução e aprovar o fluxo final;

    \item \textbf{Supervisora de Qualidade} --- responsável
    por validar a aderência do processo, as regras
    operacionais, as necessidades do setor, os impactos na
    rotina da equipe e as prioridades do fluxo;

    \item \textbf{Analista de Qualidade} --- responsável por
    descrever o processo atual, as atividades executadas, as
    dificuldades encontradas, as regras do processo, as
    funcionalidades esperadas da automação e as melhorias
    desejadas.
\end{itemize}

A escolha desses três perfis garante que o levantamento de
requisitos abranja as perspectivas estratégica, operacional e
técnica do processo, permitindo especificar uma solução que
atenda às necessidades de todos os níveis de uso. A divisão
das perguntas do formulário conforme o nível de atuação de
cada participante busca aumentar a qualidade dos requisitos
levantados e validar o fluxo proposto em diferentes níveis
organizacionais.

\subsection{Instrumento de Coleta de Dados}

A coleta de dados foi realizada por meio de um
\textbf{formulário estruturado}, elaborado na ferramenta
\textbf{Google Forms} e aplicado individualmente a cada um dos
três participantes. O formulário foi organizado em seções,
com um bloco de perguntas comuns a todos os respondentes e
blocos específicos direcionados conforme o nível de atuação
do participante (Analista, Supervisora ou Gerente), de modo
que cada perfil respondesse apenas às questões relacionadas ao
seu nível de responsabilidade.

O formulário abordou os seguintes eixos temáticos:
funcionamento atual do processo de comunicação de problemas
entre áreas; principais dificuldades enfrentadas na ausência
de um sistema formalizado; oportunidades de melhoria; regras
de negócio e exceções do processo; atividades executadas
manualmente; campos e informações consideradas essenciais no
registro de um chamado; critérios de priorização e prazos
esperados para resolução; funcionalidades esperadas da
automação; e validação do fluxo proposto em cada nível
organizacional.

As respostas coletadas foram exportadas e analisadas por meio
de análise de conteúdo, com agrupamento em categorias
temáticas que originaram os requisitos funcionais e não
funcionais da solução, posteriormente priorizados pela
técnica MoSCoW e validados pelo orientador.

\subsection{Aspectos Éticos}

A participação no formulário foi voluntária e precedida de
uma seção de apresentação da pesquisa e de consentimento
informado, na qual o participante é esclarecido sobre os
objetivos do levantamento e concorda em responder antes de
prosseguir. As informações coletadas referem-se exclusivamente
ao processo organizacional em análise --- funcionamento do
fluxo, dificuldades, regras de negócio e expectativas quanto à
automação --- não envolvendo dados pessoais sensíveis. Os
participantes são identificados no trabalho apenas pelo
cargo, sem exposição de dados que permitam sua identificação
individual.

Como limitação, destaca-se que os resultados refletem as
características específicas da organização estudada e a
perspectiva dos três participantes selecionados, não podendo
ser generalizados automaticamente para outros contextos
organizacionais, condição inerente ao delineamento de estudo
de caso \cite{Gil2002}.

\subsection{Etapas da Pesquisa}

Em coerência com o ciclo de BPM apresentado na Base Teórica
--- identificação, modelagem, análise, redesenho, automação e
monitoramento --- a pesquisa foi estruturada em seis etapas
sequenciais, apresentadas a seguir.

\subsubsection{Etapa 1 --- Diagnóstico do Processo Atual e
Revisão de Literatura}

Análise do processo existente de comunicação de chamados
internos (processo AS-IS), conduzida em paralelo à revisão de
literatura nas bases Google Scholar, SciELO e Portal de
Periódicos CAPES, com descritores relacionados a engenharia
de processos, BPM, engenharia de requisitos, automação de
processos, plataformas \textit{low-code}, Microsoft 365,
Microsoft List, Power Automate, Power BI e SLA. O objetivo é
construir a fundamentação teórica que embasa o redesenho e o
desenvolvimento da solução proposta.

\subsubsection{Etapa 2 --- Levantamento de Informações via
Google Forms}

Aplicação do formulário estruturado no Google Forms aos três
participantes da área de Qualidade (Analista, Supervisora e
Gerente), com perguntas direcionadas conforme o nível de
atuação de cada perfil. As respostas coletadas foram
analisadas e consolidadas em requisitos funcionais e não
funcionais, priorizados pela técnica MoSCoW (\textit{Must
Have, Should Have, Could Have} e \textit{Won't Have}), e
validados pelo orientador.

\subsubsection{Etapa 3 --- Especificação dos Requisitos}

Transformação das necessidades relatadas pelos participantes
em requisitos funcionais, requisitos não funcionais e regras
de negócio, seguindo a lógica necessidade~$\rightarrow$~análise~
$\rightarrow$~requisito discutida na Base Teórica, com base
nos conceitos de Engenharia de Requisitos.

\subsubsection{Etapa 4 --- Modelagem do Processo Futuro
(TO-BE)}

Redesenho do fluxo de gestão de chamados a partir dos
requisitos especificados, definindo os estados possíveis do
chamado, os critérios de prioridade, o tempo de SLA associado
a cada prioridade e as responsabilidades de cada ator do
processo, segundo a disciplina do BPM.

\subsubsection{Etapa 5 --- Desenvolvimento e Implantação do
Fluxo}

Com base no processo redesenhado, foi desenvolvido o site no
SharePoint responsável por centralizar o acesso à abertura dos
chamados, com os campos, status e critérios de prioridade
definidos na etapa anterior. Em paralelo, foi desenvolvido o
fluxo de automação no Power Automate (geração de protocolo,
notificações e contagem de SLA) e estruturada a Microsoft
List como repositório centralizado. O fluxo foi implantado em
ambiente de testes para verificação do funcionamento
integrado dos componentes, e consolidados os indicadores no
painel gerencial do Power BI (volume de chamados, tempo médio
de resolução, percentual dentro do SLA e distribuição por
área e por prioridade).

\subsubsection{Etapa 6 --- Validação e Consolidação dos
Resultados}

Acompanhamento da validação do processo em uso junto aos
participantes envolvidos, com a consolidação dos dados e
indicadores gerados pelo Power BI. Os resultados são
analisados e confrontados com a hipótese da pesquisa,
discutindo-se à luz da fundamentação teórica a aderência da
proposta às necessidades identificadas nos três níveis
organizacionais consultados.

\begin{table}[H]
\caption{Síntese das etapas da pesquisa}
\label{tab:etapas}
\centering
\begin{tabular}{p{0.5cm} p{3.6cm} p{3.6cm} p{4.4cm}}
\hline
\textbf{N} & \textbf{Etapa} & \textbf{Fonte/Método} &
\textbf{Produto Gerado} \\
\hline
1 & Diagnóstico do processo atual e revisão de literatura &
Análise do processo (AS-IS) e bases científicas &
Fundamentação teórica e caracterização do problema \\
\hline
2 & Levantamento de informações & Formulário Google Forms
com 3 participantes & Necessidades organizacionais
consolidadas \\
\hline
3 & Especificação dos requisitos & Análise de conteúdo e
técnica MoSCoW & Requisitos funcionais e não funcionais \\
\hline
4 & Modelagem do processo futuro (TO-BE) & Redesenho segundo
o ciclo de BPM & Fluxo, estados e regras de negócio
definidos \\
\hline
5 & Desenvolvimento e implantação & SharePoint + Power Apps +
Power Automate + Microsoft List + Power BI & Fluxo implantado
e dashboard lançado \\
\hline
6 & Validação e consolidação dos resultados & Análise
qualitativa junto aos participantes & Resultados e conclusões \\
\hline
\end{tabular}
\raggedright
\par
{\fontsize{10}{12}\selectfont Fonte: Elaborado pelo autor}
\end{table}

\subsection{Cronograma}

O cronograma da pesquisa está estruturado em cinco fases
mensais, distribuídas entre julho e novembro do ano de
execução, conforme descrito na Tabela~\ref{tab:cronograma}.
As Etapas 1 a 4 da pesquisa (diagnóstico, levantamento,
especificação e modelagem) concentram-se no mês de julho; as
Etapas 5 e 6 (desenvolvimento/implantação e
validação/consolidação) correspondem, respectivamente, aos
meses de agosto/setembro e outubro/novembro.

\begin{table}[H]
\caption{Cronograma de execução da pesquisa}
\label{tab:cronograma}
\centering
\begin{tabular}{p{1.5cm} p{3.5cm} p{7cm}}
\hline
\textbf{Mês} & \textbf{Fase} & \textbf{Atividades} \\
\hline
Julho &
Diagnóstico, levantamento, especificação e modelagem &
Análise do processo atual (AS-IS); elaboração e validação do
formulário no Google Forms; aplicação aos três participantes
(Gerente, Supervisora e Analista de Qualidade); análise de
conteúdo das respostas; especificação dos requisitos
funcionais e não funcionais; redesenho do processo futuro
(TO-BE); validação pelo orientador. \\
\hline
Agosto &
Desenvolvimento e implantação do fluxo &
Desenvolvimento do site no SharePoint como ponto de entrada
dos chamados; configuração dos campos obrigatórios, status e
critérios de prioridade; desenvolvimento do fluxo de
automação no Power Automate (geração de protocolo,
notificações e contagem de SLA); estruturação da Microsoft
List; implantação do fluxo em ambiente de testes. \\
\hline
Setembro &
Consolidação, indicadores e lançamento &
Consolidação da integração entre o site no SharePoint e os
demais componentes do ecossistema Microsoft 365; criação dos
indicadores e construção do painel gerencial no Power BI;
lançamento do dashboard do Power BI na primeira semana do
mês; lançamento e início da validação do novo modelo do
processo junto aos participantes. \\
\hline
Outubro &
Validação e consolidação dos resultados &
Acompanhamento da validação do processo em uso; consolidação
dos dados e indicadores gerados; análise preliminar dos
resultados frente à hipótese da pesquisa. \\
\hline
Novembro &
Validação e consolidação dos resultados (continuação) &
Consolidação final dos resultados e da documentação técnica;
redação dos resultados e discussão; revisão final do trabalho
e preparação para a apresentação. \\
\hline
\end{tabular}
\raggedright
\par
{\fontsize{10}{12}\selectfont Fonte: Elaborado pelo autor}
\end{table}

\section{Levantamento de Requisitos}

Em coerência com a Engenharia de Requisitos discutida na Base
Teórica, o levantamento apresentado nesta seção corresponde
às etapas de \textit{Análise} do ciclo de BPM: a partir do
diagnóstico do processo atual, busca-se identificar os
gargalos operacionais e transformá-los em requisitos capazes
de orientar o redesenho do processo, apresentado na seção
seguinte. O levantamento de requisitos foi realizado a partir
das respostas do formulário estruturado aplicado aos três
participantes definidos na metodologia: uma Analista de
Qualidade e Clientes, uma Supervisora de Qualidade e um
Gerente da Área. Os três participantes registraram
concordância voluntária com a pesquisa. As perguntas foram
segmentadas conforme o nível de atuação, de modo que os
campos sem resposta nos blocos específicos não representam
omissão, mas sim questões direcionadas a outro perfil
organizacional.

A análise das respostas foi conduzida por agrupamento
temático, considerando cinco eixos: processo atual,
dificuldades e riscos, regras de negócio, funcionalidades
esperadas e benefícios gerenciais. A partir desses eixos,
as necessidades relatadas foram convertidas em requisitos
funcionais, requisitos não funcionais e regras de negócio.
Para a priorização foi utilizada a técnica MoSCoW,
considerando como \textit{Must Have} os requisitos
necessários para que o fluxo opere com rastreabilidade e
governança mínimas, como \textit{Should Have} aqueles que
aumentam a qualidade operacional sem impedir a operação
básica, e como \textit{Could Have} as funcionalidades de
evolução que podem ser incorporadas em ciclos posteriores.

\subsection{Diagnóstico do Processo Atual}

A resposta da Analista de Qualidade e Clientes permitiu
caracterizar o processo operacional utilizado antes da
solução proposta. Os chamados internos eram registrados no
SAP Service, ferramenta também utilizada como CRM para
atendimentos e reclamações de clientes externos. Para
viabilizar esse registro, as ocorrências internas eram
associadas a um cadastro denominado ``Experiência do
Cliente'', no qual era criado um ticket e realizado o
direcionamento à área responsável. Paralelamente ao registro
no sistema, a chegada do chamado era comunicada manualmente
aos responsáveis por meio de um grupo de WhatsApp.

Esse modelo apresentava duas fragilidades centrais. A
primeira estava relacionada à mistura, em um mesmo ambiente,
de chamados internos e manifestações externas de clientes,
o que poderia distorcer a leitura de volumes e dificultar a
governança das ocorrências internas. A segunda estava
relacionada à dependência de atividades manuais para
comunicação e acompanhamento, especialmente quanto à
obtenção de retorno das áreas responsáveis.

A Analista também diferenciou os dois tipos de chamado
utilizados no processo. O chamado interno corresponde a uma
problemática organizacional que necessita de análise e
tratativa por outra área. O chamado massivo corresponde a
uma ocorrência de mesmo problema capaz de impactar três ou
mais clientes, normalmente associada a falhas técnicas ou
sistêmicas com maior potencial de impacto. Essa distinção
passou a ser considerada na modelagem do processo e nas
regras de priorização.

Na perspectiva da Supervisora de Qualidade, a proposta de
um ambiente específico para registro e acompanhamento foi
considerada aderente às necessidades operacionais. A
respondente destacou como benefícios esperados a maior
agilidade no direcionamento, a rastreabilidade das
tratativas, a simplificação do processo em comparação ao
SAP Service, a notificação imediata da área responsável e a
possibilidade de consolidação das informações para análise
gerencial. Como pontos de atenção, foram identificadas a
necessidade de registrar atualizações ou justificativas
quando a solução depender de fatores externos e a
necessidade de permitir a transferência do chamado quando o
responsável inicialmente definido não for o correto.

Na perspectiva gerencial, foram confirmadas como principais
fricções a dificuldade de governança dos chamados internos,
a falta de padronização dos registros e do direcionamento e
a dificuldade das áreas executoras em visualizar suas
pendências. O Gerente da Área avaliou que o fluxo proposto
possui aderência aos objetivos estratégicos da gerência e
da organização, destacando como benefícios esperados a
rastreabilidade, a padronização, a disponibilidade de
indicadores gerenciais e a redução do tempo de resolução.
O fluxo foi aprovado, em princípio, sem ressalvas.

A Tabela~\ref{tab:sintese-levantamento} sintetiza as
contribuições de cada perfil para a especificação da
solução.

\begin{table}[H]
\caption{Síntese do levantamento de requisitos por perfil}
\label{tab:sintese-levantamento}
\centering
\begin{tabular}{p{2.4cm} p{5cm} p{5cm}}
\hline
\textbf{Perfil} & \textbf{Principais evidências} &
\textbf{Implicações para a solução} \\
\hline
Analista de Qualidade e Clientes &
Uso do SAP Service para ocorrências internas e externas;
comunicação complementar por WhatsApp; dificuldade de
retorno; necessidade de protocolo, evidências, regras,
causa raiz, solução, SLA e indicadores. &
Criar canal próprio; padronizar a abertura; gerar protocolo;
automatizar notificações; registrar causa raiz e ação;
separar chamados internos de manifestações externas. \\
\hline
Supervisora de Qualidade &
Necessidade de classificação e solução por prazo;
rastreabilidade; transferência de responsável; registro de
atualizações e justificativas quando o SLA não puder ser
cumprido. &
Definir papéis de classificação; controlar mudança de
responsável; manter histórico; prever mecanismo de
justificativa e acompanhamento; estruturar regras de SLA. \\
\hline
Gerente da Área &
Dificuldade de governança, padronização e acesso às
pendências; necessidade de indicadores; aprovação do fluxo
proposto. &
Centralizar os registros; disponibilizar visão gerencial;
priorizar volume, tempo médio de resolução e distribuição
por área; garantir aderência ao processo organizacional. \\
\hline
\end{tabular}
\raggedright
\par
{\fontsize{10}{12}\selectfont Fonte: Elaborado pelo autor
com base na pesquisa de campo}
\end{table}

\subsection{Requisitos Funcionais}

Os requisitos funcionais representam os serviços que a
solução deve disponibilizar aos usuários e às áreas
envolvidas no processo. A Tabela~\ref{tab:requisitos-funcionais}
apresenta os requisitos consolidados a partir das respostas
dos participantes.

\begin{table}[H]
\caption{Requisitos funcionais da solução}
\label{tab:requisitos-funcionais}
\centering
\begin{tabular}{p{1.0cm} p{8.7cm} p{2.3cm}}
\hline
\textbf{ID} & \textbf{Requisito funcional} &
\textbf{Prioridade} \\
\hline
RF01 & Disponibilizar um canal único e específico para
abertura de chamados internos, separado do ambiente de
reclamações externas de clientes. & Must Have \\
\hline
RF02 & Permitir que usuários autorizados de diferentes
áreas registrem novos chamados internos. & Must Have \\
\hline
RF03 & Permitir a identificação do tipo de chamado,
diferenciando, no mínimo, chamados Internos e Massivos. &
Must Have \\
\hline
RF04 & Registrar de forma padronizada as informações
necessárias à abertura, incluindo empresa, tipo, assunto,
descrição do problema e evidências, além de dados
operacionais aplicáveis ao caso. & Must Have \\
\hline
RF05 & Gerar automaticamente um número único de protocolo
para cada chamado registrado. & Must Have \\
\hline
RF06 & Confirmar a abertura ao solicitante e informar o
protocolo e o link de acesso ao registro. & Must Have \\
\hline
RF07 & Notificar a equipe de Qualidade/SAC sobre novos
chamados que aguardam classificação e direcionamento. &
Must Have \\
\hline
RF08 & Permitir à Qualidade/SAC definir prioridade, área
responsável, tipologia, responsável pela tratativa e status
do chamado. & Must Have \\
\hline
RF09 & Notificar automaticamente o responsável quando um
chamado for atribuído para sua tratativa. & Must Have \\
\hline
RF10 & Permitir a transferência do chamado para outro
responsável, preservando a rastreabilidade da alteração. &
Must Have \\
\hline
RF11 & Permitir o registro de causa raiz e ação corretiva
pela área responsável pela tratativa. & Must Have \\
\hline
RF12 & Encaminhar o chamado tratado para análise da
Qualidade/SAC antes do encerramento definitivo. &
Must Have \\
\hline
RF13 & Permitir à Qualidade/SAC concluir ou cancelar o
chamado, registrando o motivo do cancelamento quando
aplicável. & Must Have \\
\hline
RF14 & Notificar o solicitante e os responsáveis quando o
chamado for concluído ou cancelado. & Must Have \\
\hline
RF15 & Controlar os prazos de classificação e solução de
acordo com o tipo e a prioridade do chamado, possibilitando
o acompanhamento do SLA. & Must Have \\
\hline
RF16 & Permitir o registro de comentários, atualizações ou
justificativas para chamados cuja tratativa dependa de
fatores externos ou ultrapasse o prazo inicialmente
definido. & Should Have \\
\hline
RF17 & Permitir o registro da quantidade de protocolos
relacionados quando a ocorrência afetar múltiplos clientes
ou registros. & Should Have \\
\hline
RF18 & Disponibilizar categorias ou tipologias reutilizáveis
para apoiar a classificação e a análise de recorrências. &
Should Have \\
\hline
RF19 & Disponibilizar indicadores gerenciais de volume,
tempo médio de resolução, distribuição por área, tipos de
problema, chamados tratados e cumprimento de SLA. &
Must Have \\
\hline
\end{tabular}
\raggedright
\par
{\fontsize{10}{12}\selectfont Fonte: Elaborado pelo autor
com base na pesquisa de campo}
\end{table}

\subsection{Requisitos Não Funcionais}

Além das funcionalidades, o levantamento evidenciou
características de qualidade necessárias para que a solução
seja incorporada à rotina organizacional. Esses requisitos
estão apresentados na Tabela~\ref{tab:requisitos-nao-funcionais}.

\begin{table}[H]
\caption{Requisitos não funcionais da solução}
\label{tab:requisitos-nao-funcionais}
\centering
\begin{tabular}{p{1.1cm} p{8.5cm} p{2.4cm}}
\hline
\textbf{ID} & \textbf{Requisito não funcional} &
\textbf{Prioridade} \\
\hline
RNF01 & A solução deve centralizar os registros em uma fonte
única e permitir rastreabilidade desde a abertura até o
encerramento. & Must Have \\
\hline
RNF02 & A interface de abertura deve ser simples, objetiva e
adequada ao uso em ambiente corporativo, reduzindo o número
de etapas necessárias para registrar uma ocorrência. &
Must Have \\
\hline
RNF03 & A solução deve utilizar, prioritariamente, os
recursos do ecossistema Microsoft 365 já adotado pela
organização, favorecendo integração e manutenção. &
Must Have \\
\hline
RNF04 & O acesso aos registros e às ações do processo deve
respeitar as permissões corporativas e os diferentes papéis
de solicitante, Qualidade/SAC e área responsável. &
Must Have \\
\hline
RNF05 & As alterações relevantes de status, responsável,
tratativa e encerramento devem permanecer auditáveis no
histórico do registro. & Must Have \\
\hline
RNF06 & As automações devem evitar notificações duplicadas e
reprocessamentos causados pelas próprias atualizações
realizadas no item. & Must Have \\
\hline
RNF07 & A interface incorporada ao SharePoint deve adaptar-se
ao espaço disponível na página, preservando legibilidade e
usabilidade. & Should Have \\
\hline
\end{tabular}
\raggedright
\par
{\fontsize{10}{12}\selectfont Fonte: Elaborado pelo autor}
\end{table}

\subsection{Regras de Negócio}

As respostas da pesquisa permitiram definir regras
operacionais que orientam o comportamento do fluxo e a
distribuição de responsabilidades entre os participantes.
A Tabela~\ref{tab:regras-negocio} apresenta as principais
regras identificadas.

\begin{table}[H]
\caption{Regras de negócio identificadas}
\label{tab:regras-negocio}
\centering
\begin{tabular}{p{1.0cm} p{10.9cm}}
\hline
\textbf{ID} & \textbf{Regra de negócio} \\
\hline
RN01 & Todo chamado deve ser registrado assim que a
problemática for identificada, garantindo rastreabilidade
desde o início do processo. \\
\hline
RN02 & A Qualidade/SAC é responsável pela classificação e
pelo direcionamento inicial do chamado para a área
responsável pela tratativa. \\
\hline
RN03 & Chamados Massivos devem possuir tratamento mais
restritivo em razão do potencial impacto sobre maior
quantidade de clientes e indicadores. \\
\hline
RN04 & Na pesquisa de campo, a classificação proposta para
Chamados Massivos foi Alta e para Chamados Internos foi
Média. \\
\hline
RN05 & A classificação inicial de um Chamado Massivo deve
ocorrer em até 30 minutos após a abertura; para Chamados
Internos, em até um dia após a abertura. \\
\hline
RN06 & A solução de um Chamado Massivo deve ocorrer em até
24 horas após a classificação; para Chamados Internos, em
até três dias após a classificação. \\
\hline
RN07 & A área responsável deve registrar, no mínimo, a causa
raiz identificada e a ação corretiva adotada antes de
submeter a tratativa para análise da Qualidade/SAC. \\
\hline
RN08 & A conclusão definitiva do chamado deve ser realizada
pela Qualidade/SAC após análise da tratativa registrada pela
área responsável. \\
\hline
RN09 & Quando houver transferência de responsável, a
alteração deve preservar a rastreabilidade do chamado e
notificar o novo responsável. \\
\hline
RN10 & Em situações extraordinárias de impacto imediato,
como indisponibilidade generalizada de sistemas, o processo
deve permitir tratamento prioritário sem impedir o registro
formal da ocorrência. \\
\hline
\end{tabular}
\raggedright
\par
{\fontsize{10}{12}\selectfont Fonte: Elaborado pelo autor
com base na pesquisa de campo}
\end{table}

\subsection{Decisões de Modelagem Derivadas dos Requisitos}

Durante a transformação dos requisitos em solução, algumas
decisões de modelagem foram adotadas para reduzir erros de
direcionamento e separar responsabilidades. Embora a
Analista tenha indicado ``Área responsável'' entre as
informações desejáveis na abertura, a mesma respondente
também solicitou que a área de Clientes/Qualidade atuasse
como filtro interno para identificar o destino correto.
A Supervisora, por sua vez, apontou a necessidade de
transferência quando o direcionamento inicial estivesse
incorreto. Em razão dessa convergência, a definição da área
responsável, da prioridade, da tipologia e do responsável
pela tratativa foi posicionada na etapa de classificação,
sob responsabilidade da Qualidade/SAC, e não no formulário
inicial do solicitante.

Também foi estabelecida uma etapa intermediária entre a
tratativa e o encerramento. Após o preenchimento da causa
raiz e da ação corretiva, o chamado deve retornar para
análise da Qualidade/SAC. Essa decisão atende à necessidade
de obrigatoriedade de posicionamento, análise de causa e
solução indicada no levantamento e impede que a própria área
responsável encerre diretamente um chamado sem validação.

Por fim, o levantamento demonstrou que a automação não deve
ser compreendida apenas como mecanismo de registro. A
necessidade de governança, rastreabilidade e indicadores
foi recorrente nos três níveis consultados, tornando a
centralização dos dados e a manutenção de um histórico
estruturado elementos centrais da solução.

\section{Levantamento de Requisitos}

Em coerência com a Engenharia de Requisitos discutida na Base
Teórica, o levantamento apresentado nesta seção corresponde
às etapas de \textit{Análise} do ciclo de BPM: a partir do
diagnóstico do processo atual, busca-se identificar os
gargalos operacionais e transformá-los em requisitos capazes
de orientar o redesenho do processo, apresentado na seção
seguinte. O levantamento de requisitos foi realizado a partir
das respostas do formulário estruturado aplicado aos três
participantes definidos na metodologia: uma Analista de
Qualidade e Clientes, uma Supervisora de Qualidade e um
Gerente da Área. Os três participantes registraram
concordância voluntária com a pesquisa. As perguntas foram
segmentadas conforme o nível de atuação, de modo que os
campos sem resposta nos blocos específicos não representam
omissão, mas sim questões direcionadas a outro perfil
organizacional.

A análise das respostas foi conduzida por agrupamento
temático, considerando cinco eixos: processo atual,
dificuldades e riscos, regras de negócio, funcionalidades
esperadas e benefícios gerenciais. A partir desses eixos,
as necessidades relatadas foram convertidas em requisitos
funcionais, requisitos não funcionais e regras de negócio.
Para a priorização foi utilizada a técnica MoSCoW,
considerando como \textit{Must Have} os requisitos
necessários para que o fluxo opere com rastreabilidade e
governança mínimas, como \textit{Should Have} aqueles que
aumentam a qualidade operacional sem impedir a operação
básica, e como \textit{Could Have} as funcionalidades de
evolução que podem ser incorporadas em ciclos posteriores.

\subsection{Diagnóstico do Processo Atual}

A resposta da Analista de Qualidade e Clientes permitiu
caracterizar o processo operacional utilizado antes da
solução proposta. Os chamados internos eram registrados no
SAP Service, ferramenta também utilizada como CRM para
atendimentos e reclamações de clientes externos. Para
viabilizar esse registro, as ocorrências internas eram
associadas a um cadastro denominado ``Experiência do
Cliente'', no qual era criado um ticket e realizado o
direcionamento à área responsável. Paralelamente ao registro
no sistema, a chegada do chamado era comunicada manualmente
aos responsáveis por meio de um grupo de WhatsApp.

Esse modelo apresentava duas fragilidades centrais. A
primeira estava relacionada à mistura, em um mesmo ambiente,
de chamados internos e manifestações externas de clientes,
o que poderia distorcer a leitura de volumes e dificultar a
governança das ocorrências internas. A segunda estava
relacionada à dependência de atividades manuais para
comunicação e acompanhamento, especialmente quanto à
obtenção de retorno das áreas responsáveis.

A Analista também diferenciou os dois tipos de chamado
utilizados no processo. O chamado interno corresponde a uma
problemática organizacional que necessita de análise e
tratativa por outra área. O chamado massivo corresponde a
uma ocorrência de mesmo problema capaz de impactar três ou
mais clientes, normalmente associada a falhas técnicas ou
sistêmicas com maior potencial de impacto. Essa distinção
passou a ser considerada na modelagem do processo e nas
regras de priorização.

Na perspectiva da Supervisora de Qualidade, a proposta de
um ambiente específico para registro e acompanhamento foi
considerada aderente às necessidades operacionais. A
respondente destacou como benefícios esperados a maior
agilidade no direcionamento, a rastreabilidade das
tratativas, a simplificação do processo em comparação ao
SAP Service, a notificação imediata da área responsável e a
possibilidade de consolidação das informações para análise
gerencial. Como pontos de atenção, foram identificadas a
necessidade de registrar atualizações ou justificativas
quando a solução depender de fatores externos e a
necessidade de permitir a transferência do chamado quando o
responsável inicialmente definido não for o correto.

Na perspectiva gerencial, foram confirmadas como principais
fricções a dificuldade de governança dos chamados internos,
a falta de padronização dos registros e do direcionamento e
a dificuldade das áreas executoras em visualizar suas
pendências. O Gerente da Área avaliou que o fluxo proposto
possui aderência aos objetivos estratégicos da gerência e
da organização, destacando como benefícios esperados a
rastreabilidade, a padronização, a disponibilidade de
indicadores gerenciais e a redução do tempo de resolução.
O fluxo foi aprovado, em princípio, sem ressalvas.

A Tabela~\ref{tab:sintese-levantamento} sintetiza as
contribuições de cada perfil para a especificação da
solução.

\begin{table}[H]
\caption{Síntese do levantamento de requisitos por perfil}
\label{tab:sintese-levantamento}
\centering
\begin{tabular}{p{2.4cm} p{5cm} p{5cm}}
\hline
\textbf{Perfil} & \textbf{Principais evidências} &
\textbf{Implicações para a solução} \\
\hline
Analista de Qualidade e Clientes &
Uso do SAP Service para ocorrências internas e externas;
comunicação complementar por WhatsApp; dificuldade de
retorno; necessidade de protocolo, evidências, regras,
causa raiz, solução, SLA e indicadores. &
Criar canal próprio; padronizar a abertura; gerar protocolo;
automatizar notificações; registrar causa raiz e ação;
separar chamados internos de manifestações externas. \\
\hline
Supervisora de Qualidade &
Necessidade de classificação e solução por prazo;
rastreabilidade; transferência de responsável; registro de
atualizações e justificativas quando o SLA não puder ser
cumprido. &
Definir papéis de classificação; controlar mudança de
responsável; manter histórico; prever mecanismo de
justificativa e acompanhamento; estruturar regras de SLA. \\
\hline
Gerente da Área &
Dificuldade de governança, padronização e acesso às
pendências; necessidade de indicadores; aprovação do fluxo
proposto. &
Centralizar os registros; disponibilizar visão gerencial;
priorizar volume, tempo médio de resolução e distribuição
por área; garantir aderência ao processo organizacional. \\
\hline
\end{tabular}
\raggedright
\par
{\fontsize{10}{12}\selectfont Fonte: Elaborado pelo autor
com base na pesquisa de campo}
\end{table}

\subsection{Requisitos Funcionais}

Os requisitos funcionais representam os serviços que a
solução deve disponibilizar aos usuários e às áreas
envolvidas no processo. A Tabela~\ref{tab:requisitos-funcionais}
apresenta os requisitos consolidados a partir das respostas
dos participantes.

\begin{table}[H]
\caption{Requisitos funcionais da solução}
\label{tab:requisitos-funcionais}
\centering
\begin{tabular}{p{1.0cm} p{8.7cm} p{2.3cm}}
\hline
\textbf{ID} & \textbf{Requisito funcional} &
\textbf{Prioridade} \\
\hline
RF01 & Disponibilizar um canal único e específico para
abertura de chamados internos, separado do ambiente de
reclamações externas de clientes. & Must Have \\
\hline
RF02 & Permitir que usuários autorizados de diferentes
áreas registrem novos chamados internos. & Must Have \\
\hline
RF03 & Permitir a identificação do tipo de chamado,
diferenciando, no mínimo, chamados Internos e Massivos. &
Must Have \\
\hline
RF04 & Registrar de forma padronizada as informações
necessárias à abertura, incluindo empresa, tipo, assunto,
descrição do problema e evidências, além de dados
operacionais aplicáveis ao caso. & Must Have \\
\hline
RF05 & Gerar automaticamente um número único de protocolo
para cada chamado registrado. & Must Have \\
\hline
RF06 & Confirmar a abertura ao solicitante e informar o
protocolo e o link de acesso ao registro. & Must Have \\
\hline
RF07 & Notificar a equipe de Qualidade/SAC sobre novos
chamados que aguardam classificação e direcionamento. &
Must Have \\
\hline
RF08 & Permitir à Qualidade/SAC definir prioridade, área
responsável, tipologia, responsável pela tratativa e status
do chamado. & Must Have \\
\hline
RF09 & Notificar automaticamente o responsável quando um
chamado for atribuído para sua tratativa. & Must Have \\
\hline
RF10 & Permitir a transferência do chamado para outro
responsável, preservando a rastreabilidade da alteração. &
Must Have \\
\hline
RF11 & Permitir o registro de causa raiz e ação corretiva
pela área responsável pela tratativa. & Must Have \\
\hline
RF12 & Encaminhar o chamado tratado para análise da
Qualidade/SAC antes do encerramento definitivo. &
Must Have \\
\hline
RF13 & Permitir à Qualidade/SAC concluir ou cancelar o
chamado, registrando o motivo do cancelamento quando
aplicável. & Must Have \\
\hline
RF14 & Notificar o solicitante e os responsáveis quando o
chamado for concluído ou cancelado. & Must Have \\
\hline
RF15 & Controlar os prazos de classificação e solução de
acordo com o tipo e a prioridade do chamado, possibilitando
o acompanhamento do SLA. & Must Have \\
\hline
RF16 & Permitir o registro de comentários, atualizações ou
justificativas para chamados cuja tratativa dependa de
fatores externos ou ultrapasse o prazo inicialmente
definido. & Should Have \\
\hline
RF17 & Permitir o registro da quantidade de protocolos
relacionados quando a ocorrência afetar múltiplos clientes
ou registros. & Should Have \\
\hline
RF18 & Disponibilizar categorias ou tipologias reutilizáveis
para apoiar a classificação e a análise de recorrências. &
Should Have \\
\hline
RF19 & Disponibilizar indicadores gerenciais de volume,
tempo médio de resolução, distribuição por área, tipos de
problema, chamados tratados e cumprimento de SLA. &
Must Have \\
\hline
\end{tabular}
\raggedright
\par
{\fontsize{10}{12}\selectfont Fonte: Elaborado pelo autor
com base na pesquisa de campo}
\end{table}

\subsection{Requisitos Não Funcionais}

Além das funcionalidades, o levantamento evidenciou
características de qualidade necessárias para que a solução
seja incorporada à rotina organizacional. Esses requisitos
estão apresentados na Tabela~\ref{tab:requisitos-nao-funcionais}.

\begin{table}[H]
\caption{Requisitos não funcionais da solução}
\label{tab:requisitos-nao-funcionais}
\centering
\begin{tabular}{p{1.1cm} p{8.5cm} p{2.4cm}}
\hline
\textbf{ID} & \textbf{Requisito não funcional} &
\textbf{Prioridade} \\
\hline
RNF01 & A solução deve centralizar os registros em uma fonte
única e permitir rastreabilidade desde a abertura até o
encerramento. & Must Have \\
\hline
RNF02 & A interface de abertura deve ser simples, objetiva e
adequada ao uso em ambiente corporativo, reduzindo o número
de etapas necessárias para registrar uma ocorrência. &
Must Have \\
\hline
RNF03 & A solução deve utilizar, prioritariamente, os
recursos do ecossistema Microsoft 365 já adotado pela
organização, favorecendo integração e manutenção. &
Must Have \\
\hline
RNF04 & O acesso aos registros e às ações do processo deve
respeitar as permissões corporativas e os diferentes papéis
de solicitante, Qualidade/SAC e área responsável. &
Must Have \\
\hline
RNF05 & As alterações relevantes de status, responsável,
tratativa e encerramento devem permanecer auditáveis no
histórico do registro. & Must Have \\
\hline
RNF06 & As automações devem evitar notificações duplicadas e
reprocessamentos causados pelas próprias atualizações
realizadas no item. & Must Have \\
\hline
RNF07 & A interface incorporada ao SharePoint deve adaptar-se
ao espaço disponível na página, preservando legibilidade e
usabilidade. & Should Have \\
\hline
\end{tabular}
\raggedright
\par
{\fontsize{10}{12}\selectfont Fonte: Elaborado pelo autor}
\end{table}

\subsection{Regras de Negócio}

As respostas da pesquisa permitiram definir regras
operacionais que orientam o comportamento do fluxo e a
distribuição de responsabilidades entre os participantes.
A Tabela~\ref{tab:regras-negocio} apresenta as principais
regras identificadas.

\begin{table}[H]
\caption{Regras de negócio identificadas}
\label{tab:regras-negocio}
\centering
\begin{tabular}{p{1.0cm} p{10.9cm}}
\hline
\textbf{ID} & \textbf{Regra de negócio} \\
\hline
RN01 & Todo chamado deve ser registrado assim que a
problemática for identificada, garantindo rastreabilidade
desde o início do processo. \\
\hline
RN02 & A Qualidade/SAC é responsável pela classificação e
pelo direcionamento inicial do chamado para a área
responsável pela tratativa. \\
\hline
RN03 & Chamados Massivos devem possuir tratamento mais
restritivo em razão do potencial impacto sobre maior
quantidade de clientes e indicadores. \\
\hline
RN04 & Na pesquisa de campo, a classificação proposta para
Chamados Massivos foi Alta e para Chamados Internos foi
Média. \\
\hline
RN05 & A classificação inicial de um Chamado Massivo deve
ocorrer em até 30 minutos após a abertura; para Chamados
Internos, em até um dia após a abertura. \\
\hline
RN06 & A solução de um Chamado Massivo deve ocorrer em até
24 horas após a classificação; para Chamados Internos, em
até três dias após a classificação. \\
\hline
RN07 & A área responsável deve registrar, no mínimo, a causa
raiz identificada e a ação corretiva adotada antes de
submeter a tratativa para análise da Qualidade/SAC. \\
\hline
RN08 & A conclusão definitiva do chamado deve ser realizada
pela Qualidade/SAC após análise da tratativa registrada pela
área responsável. \\
\hline
RN09 & Quando houver transferência de responsável, a
alteração deve preservar a rastreabilidade do chamado e
notificar o novo responsável. \\
\hline
RN10 & Em situações extraordinárias de impacto imediato,
como indisponibilidade generalizada de sistemas, o processo
deve permitir tratamento prioritário sem impedir o registro
formal da ocorrência. \\
\hline
\end{tabular}
\raggedright
\par
{\fontsize{10}{12}\selectfont Fonte: Elaborado pelo autor
com base na pesquisa de campo}
\end{table}

\subsection{Decisões de Modelagem Derivadas dos Requisitos}

Durante a transformação dos requisitos em solução, algumas
decisões de modelagem foram adotadas para reduzir erros de
direcionamento e separar responsabilidades. Embora a
Analista tenha indicado ``Área responsável'' entre as
informações desejáveis na abertura, a mesma respondente
também solicitou que a área de Clientes/Qualidade atuasse
como filtro interno para identificar o destino correto.
A Supervisora, por sua vez, apontou a necessidade de
transferência quando o direcionamento inicial estivesse
incorreto. Em razão dessa convergência, a definição da área
responsável, da prioridade, da tipologia e do responsável
pela tratativa foi posicionada na etapa de classificação,
sob responsabilidade da Qualidade/SAC, e não no formulário
inicial do solicitante.

Também foi estabelecida uma etapa intermediária entre a
tratativa e o encerramento. Após o preenchimento da causa
raiz e da ação corretiva, o chamado deve retornar para
análise da Qualidade/SAC. Essa decisão atende à necessidade
de obrigatoriedade de posicionamento, análise de causa e
solução indicada no levantamento e impede que a própria área
responsável encerre diretamente um chamado sem validação.

Por fim, o levantamento demonstrou que a automação não deve
ser compreendida apenas como mecanismo de registro. A
necessidade de governança, rastreabilidade e indicadores
foi recorrente nos três níveis consultados, tornando a
centralização dos dados e a manutenção de um histórico
estruturado elementos centrais da solução.

\section{Análise e Redesenho do Processo}

Esta seção materializa a etapa de \textit{Redesenho} do ciclo
de BPM apresentado na Base Teórica, articulando o diagnóstico
do processo atual (AS-IS), descrito na seção de Levantamento
de Requisitos, com a especificação do processo futuro
(TO-BE), que orienta o desenvolvimento da solução apresentado
na seção seguinte. O objetivo não é apenas contrastar duas
representações do fluxo, mas evidenciar, campo a campo,
decisão a decisão, a origem de cada mudança nos requisitos
levantados junto aos três perfis consultados.

\subsection{Processo Atual (AS-IS)}

O diagnóstico consolidado a partir das respostas da Analista,
da Supervisora e do Gerente caracteriza um processo
descentralizado e dependente de controles manuais. Os
chamados internos eram registrados no SAP Service --- mesma
ferramenta utilizada como CRM para atendimentos e reclamações
de clientes externos ---, misturando ocorrências internas e
manifestações externas em um único ambiente. A chegada do
chamado era comunicada paralelamente por um grupo de
WhatsApp, sem protocolo automático, sem controle de prazo e
sem visibilidade gerencial sobre volume, prioridade ou tempo
de resolução. O Quadro~\ref{quadro:as-is} sintetiza as
características centrais do processo atual.

\begin{table}[H]
\caption{Características do processo atual (AS-IS)}
\label{quadro:as-is}
\centering
\begin{tabular}{p{3.6cm} p{8.2cm}}
\hline
\textbf{Dimensão} & \textbf{Situação observada} \\
\hline
Canal de abertura & Cadastro ``Experiência do Cliente'' no
SAP Service, compartilhado com manifestações externas de
clientes. \\
\hline
Comunicação & Complementar e manual, por grupo de WhatsApp,
sem registro formal do repasse. \\
\hline
Controle de prazo & Inexistente; sem SLA definido ou
monitorado. \\
\hline
Rastreabilidade & Baixa; mistura entre chamados internos e
externos dificulta a leitura de volumes. \\
\hline
Responsabilidade & Informal; dependência de retorno manual
das áreas responsáveis. \\
\hline
Indicadores & Inexistentes; ausência de base para gestão
baseada em dados. \\
\hline
\end{tabular}
\raggedright
\par
{\fontsize{10}{12}\selectfont Fonte: Elaborado pelo autor com
base na pesquisa de campo}
\end{table}

\subsection{Processo Futuro (TO-BE)}

O processo redesenhado responde a cada uma das fragilidades
identificadas no AS-IS, substituindo controles manuais por
regras automatizadas e centralizando o registro em uma única
fonte de dados. O Quadro~\ref{quadro:to-be} apresenta a
transformação proposta, dimensão a dimensão, em correspondência
direta com o Quadro~\ref{quadro:as-is}.

\begin{table}[H]
\caption{Transformação do processo: AS-IS para TO-BE}
\label{quadro:to-be}
\centering
\begin{tabular}{p{3.0cm} p{4.9cm} p{4.9cm}}
\hline
\textbf{Dimensão} & \textbf{AS-IS} & \textbf{TO-BE} \\
\hline
Canal de abertura & Cadastro compartilhado no SAP Service &
Site no SharePoint com aplicativo Power Apps próprio para
chamados internos \\
\hline
Comunicação & Grupo de WhatsApp, sem registro & Notificações
automáticas por e-mail, vinculadas a eventos da Microsoft
List \\
\hline
Controle de prazo & Inexistente & SLA por tipo e prioridade,
com contagem automatizada \\
\hline
Rastreabilidade & Baixa, dados misturados & Protocolo único,
histórico centralizado na Microsoft List \\
\hline
Responsabilidade & Informal & Classificação e atribuição
explícita pela Qualidade/SAC, com transferência rastreável \\
\hline
Indicadores & Inexistentes & Painel gerencial no Power BI
(volume, tempo médio, cumprimento de SLA) \\
\hline
\end{tabular}
\raggedright
\par
{\fontsize{10}{12}\selectfont Fonte: Elaborado pelo autor}
\end{table}

\subsection{Decisões de Redesenho}

A transição entre os dois estados do processo não decorreu da
transposição direta do fluxo manual para um ambiente digital,
mas de decisões de redesenho fundamentadas nos requisitos
levantados. Duas decisões merecem destaque.

A primeira separa o registro de chamados internos das
manifestações externas de clientes, respondendo diretamente à
fragilidade central identificada pela Analista no processo
AS-IS. Essa separação é também uma decisão de Engenharia de
Processos: ao isolar o fluxo interno em um canal próprio, a
organização ganha capacidade de leitura de volumes e de
governança que o modelo anterior não permitia.

A segunda desloca a definição de área responsável, prioridade,
tipologia e responsável pela tratativa do formulário inicial
de abertura para uma etapa intermediária de classificação,
sob responsabilidade da Qualidade/SAC. Essa decisão resulta da
convergência entre a necessidade de filtro interno apontada
pela Analista e a necessidade de transferência de responsável
apontada pela Supervisora, e ilustra como o redesenho do
processo (TO-BE) é uma consequência direta do levantamento de
requisitos, e não uma escolha arbitrária de arquitetura de
software.

Por fim, foi estabelecida uma etapa de retorno obrigatório à
Qualidade/SAC após o preenchimento da causa raiz e da ação
corretiva pela área responsável, antes do encerramento
definitivo do chamado. Essa regra opera como mecanismo de
validação previsto pelo ciclo de BPM, impedindo que a própria
área executora encerre um chamado sem revisão, e reforça a
rastreabilidade e a governança como objetivos centrais do
redesenho.

O processo futuro (TO-BE) especificado nesta seção é a base
direta para os componentes desenvolvidos e descritos na seção
de Desenvolvimento da Solução, na qual cada estado, campo e
regra de negócio aqui definidos é implementado no ecossistema
Microsoft 365.

\section{Desenvolvimento da Solução}

Esta seção corresponde à etapa de \textit{Automação} do
ciclo de BPM apresentado na Base Teórica: o processo futuro
(TO-BE) especificado na seção de Análise e Redesenho do
Processo é aqui transformado em componentes executáveis do
ecossistema Microsoft 365, preservando, mesmo em ambiente
\textit{low-code}, as etapas de arquitetura, implementação e
validação próprias da Engenharia de Software. O
desenvolvimento foi realizado com base nos requisitos
consolidados na etapa anterior e utilizou componentes do
ecossistema Microsoft 365 já disponível na organização. A
solução resultante foi estruturada em cinco elementos
integrados: uma Central de Chamados no SharePoint, um
aplicativo Canvas no Power Apps para abertura das
solicitações, uma Microsoft List denominada
\textit{Chamados} como repositório central, dois fluxos
principais no Power Automate e uma área de acesso aos
indicadores gerenciais.

A arquitetura implementada difere parcialmente da
arquitetura inicialmente descrita na proposta. Na versão
desenvolvida, o SharePoint atua como portal de acesso,
enquanto o registro do chamado é realizado por um
aplicativo Power Apps incorporado à página. O Power Apps
grava diretamente o novo item na Microsoft List e, a partir
da criação desse item, o fluxo FL01 é disparado para gerar o
protocolo e realizar as notificações. Essa alteração
preservou o SharePoint como ponto único de entrada, porém
separou a camada de interface da camada de armazenamento,
permitindo maior controle sobre a experiência do
solicitante.

\subsection{Arquitetura da Solução Desenvolvida}

A Tabela~\ref{tab:arquitetura-desenvolvida} apresenta os
componentes efetivamente utilizados e suas responsabilidades
na versão desenvolvida.

\begin{table}[H]
\caption{Componentes da arquitetura desenvolvida}
\label{tab:arquitetura-desenvolvida}
\centering
\begin{tabular}{p{3.0cm} p{4.1cm} p{5.1cm}}
\hline
\textbf{Componente} & \textbf{Função principal} &
\textbf{Responsabilidade no fluxo} \\
\hline
SharePoint &
Portal corporativo da solução &
Centralizar o acesso para abertura de chamado, consulta de
registros e acesso ao painel gerencial. \\
\hline
Power Apps &
Interface de abertura &
Coletar os dados do solicitante, validar campos
obrigatórios, permitir anexos e gravar o novo chamado na
Microsoft List. \\
\hline
Microsoft List &
Repositório central &
Armazenar dados de abertura, classificação, responsável,
tratativa, status, datas e histórico de alterações. \\
\hline
Power Automate --- FL01 &
Automação da abertura &
Gerar protocolo, inicializar dados do chamado e notificar
Qualidade/SAC e solicitante. \\
\hline
Power Automate --- FL03 &
Automação da tratativa e encerramento &
Controlar notificações por mudança de status, transferência
de responsável, envio para análise, conclusão e
cancelamento. \\
\hline
Power BI &
Camada gerencial &
Disponibilizar a visualização dos indicadores e da base
histórica a partir do portal da Central de Chamados. \\
\hline
\end{tabular}
\raggedright
\par
{\fontsize{10}{12}\selectfont Fonte: Elaborado pelo autor}
\end{table}

O fluxo lógico da solução pode ser representado pela
sequência: solicitante acessa o SharePoint, realiza a
abertura no Power Apps, o registro é persistido na
Microsoft List, o FL01 gera protocolo e notificações, a
Qualidade/SAC classifica e direciona o chamado, a área
responsável realiza a tratativa e o FL03 controla as
mudanças subsequentes até o encerramento.

\subsection{Central de Chamados no SharePoint}

O SharePoint foi utilizado como camada de apresentação e
ponto central de acesso ao processo. A página desenvolvida
apresenta uma identificação institucional da Central de
Chamados Internos e disponibiliza acessos para abertura de
novo chamado, consulta dos chamados registrados e acesso ao
dashboard gerencial. Essa organização concentra em um único
ambiente as funções que, no processo anterior, estavam
distribuídas entre SAP Service, mensagens e comunicações
paralelas.

A utilização do SharePoint como portal também permite que o
processo seja incorporado ao ambiente corporativo já
utilizado pelos usuários, reduzindo a necessidade de
acesso a uma aplicação externa. O formulário de abertura
foi incorporado à página por meio de um aplicativo Power
Apps, enquanto a Microsoft List permanece responsável pelo
armazenamento e pelas etapas internas de classificação,
tratativa e encerramento.

\subsection{Aplicativo Power Apps para Abertura do Chamado}

A primeira etapa do formulário originalmente prevista
diretamente na Microsoft List foi substituída por um
aplicativo Canvas desenvolvido no Power Apps. A mudança foi
realizada com o objetivo de oferecer uma interface mais
simples ao solicitante e separar os campos de abertura dos
campos internos de classificação e tratamento.

O aplicativo foi integrado diretamente à lista
\textit{Chamados} e configurado para criação de novos
registros. Na versão analisada, o formulário de abertura
contém os campos Empresa, Tipo de chamado, Assunto do
chamado, Problema, Linha, Veículo e Anexos. Empresa, Tipo
de chamado, Assunto do chamado e Problema foram
configurados como obrigatórios. O campo Problema utiliza
entrada de texto em múltiplas linhas, possibilitando uma
descrição detalhada da ocorrência, enquanto o controle de
Anexos permite incluir evidências que apoiem a análise.

Os campos de escolha são carregados a partir das opções
cadastradas na própria Microsoft List, evitando a
manutenção de listas paralelas no aplicativo. O botão
\textit{Criar Chamado} utiliza a validação do formulário
antes do envio, enquanto o botão \textit{Limpar Formulário}
reinicializa os campos para uma nova entrada.

Após o salvamento bem-sucedido, o aplicativo altera seu
estado visual e apresenta ao usuário uma confirmação de que
o chamado foi aberto e está aguardando análise da
Qualidade/SAC. A mesma tela disponibiliza o botão
\textit{Novo Chamado}, que limpa o formulário, retorna ao
modo de criação e permite novo registro. Em caso de falha
no envio, uma mensagem de erro é apresentada ao usuário.

A interface foi ajustada para uso incorporado em página do
SharePoint, utilizando contêineres responsivos e dimensões
baseadas no espaço disponível. Essa adequação buscou
preservar a legibilidade do formulário quando executado
dentro do \textit{web part} do Power Apps.

\subsection{Estrutura da Microsoft List}

A Microsoft List denominada \textit{Chamados} constitui o
repositório central da solução. A estrutura separa os dados
informados pelo solicitante dos dados preenchidos durante o
processamento interno.

Na abertura são armazenados, entre outros, Empresa, Tipo de
chamado, Assunto, Problema, Veículo, Linha, anexos e o
usuário responsável pela criação do item. Nas etapas
posteriores são utilizados os campos Prioridade, Área
Responsável, Tipologia, Responsável pela tratativa, Causa
raiz, Ação corretiva, Data da Tratativa, Status do Chamado,
dados de transferência, Data de Fechamento e controle de
notificação de fechamento.

O registro centralizado permite que todas as etapas
utilizem o mesmo item como referência, evitando a criação
de bases paralelas e mantendo o histórico de alterações
associado ao protocolo. A estrutura também sustenta as
automações do Power Automate, que consultam alterações de
campos específicos para identificar quando uma nova ação
deve ser executada.

\subsection{Fluxo FL01 --- Abertura de Chamados Internos}

O fluxo FL01 é iniciado automaticamente quando um novo item
é criado na Microsoft List. Sua primeira ação calcula o
protocolo utilizando o ano corrente e o identificador
numérico do item. O padrão implementado é
\texttt{C-AAAA-NNNN}. Assim, um item de identificador 14
criado em 2026, por exemplo, recebe o protocolo
\texttt{C-2026-0014}.

Após a geração do protocolo, o próprio item é atualizado.
Na versão exportada do fluxo, o status inicial é definido
como \textit{Aberto} e a prioridade é inicialmente
registrada como \textit{Média}. Em seguida são executadas
duas notificações distintas.

A primeira notificação é destinada à equipe de
Qualidade/SAC e informa que um novo chamado aguarda
classificação. O e-mail apresenta protocolo, solicitante,
empresa, tipo, assunto, linha, veículo, status e descrição
do problema, além de um link direto para o item. O objetivo
é iniciar a etapa de classificação e direcionamento.

A segunda notificação é destinada ao próprio solicitante e
funciona como confirmação formal de abertura. O e-mail
informa o número do protocolo, resume os dados registrados,
indica que a solicitação aguarda análise da Qualidade/SAC e
fornece o link para acompanhamento do chamado. Com isso, a
confirmação de abertura deixa de depender de comunicação
manual.

\subsection{Classificação e Direcionamento pela
Qualidade/SAC}

Após a abertura, a Qualidade/SAC realiza a classificação
diretamente no registro da Microsoft List. Nessa etapa são
definidos os campos Prioridade, Área Responsável, Tipologia
e Responsável pela tratativa. O status do chamado também é
controlado pela Qualidade/SAC.

A decisão de manter esses campos fora do formulário do
solicitante reduz o risco de encaminhamento incorreto e
formaliza a Qualidade/SAC como instância de triagem. Quando
o chamado é validado e direcionado para tratamento, o
status é alterado para \textit{Em tratativa}, iniciando a
execução das regras do FL03 associadas à atribuição do
responsável.

A estrutura da lista também contempla o campo Motivo do
Cancelamento. Quando a Qualidade/SAC decide não prosseguir
com a solicitação e altera o status para
\textit{Cancelado}, o motivo registrado passa a compor a
rastreabilidade do encerramento sem tratativa.

\subsection{Fluxo FL03 --- Tratativa e Encerramento}

O FL03 é disparado quando um item da lista é modificado.
Inicialmente, o fluxo obtém quais propriedades foram
alteradas e recupera o estado atual do item. Em seguida,
uma estrutura de decisão do tipo \textit{Switch} avalia o
status do chamado e direciona o processamento para o ramo
correspondente.

No estado \textit{Em tratativa}, o fluxo verifica se houve
mudança do status para identificar uma nova atribuição.
Quando essa condição é atendida, o responsável pela
tratativa recebe uma notificação por e-mail contendo os
dados necessários e o link para acesso ao chamado.

Ainda no estado \textit{Em tratativa}, o fluxo verifica se
o campo de responsável foi alterado. Quando ocorre uma
transferência, o registro é atualizado com informações de
controle da transferência e o novo responsável é
notificado. Essa regra foi incorporada após a necessidade
identificada no levantamento de requisitos de redistribuir
chamados sem perda de rastreabilidade.

Quando Causa raiz e Ação corretiva são preenchidas e as
respectivas propriedades são identificadas como alteradas,
o FL03 atualiza automaticamente o status para
\textit{Em Análise} e envia uma notificação à
Qualidade/SAC. A área responsável, portanto, registra a
tratativa, mas não realiza o encerramento definitivo.

A análise final é responsabilidade da Qualidade/SAC. Após
avaliar a tratativa, a equipe pode alterar manualmente o
status para \textit{Concluído} ou, quando aplicável,
\textit{Cancelado}. No caso de conclusão, o FL03 notifica o
solicitante e os responsáveis envolvidos e registra a data
de fechamento.

Para impedir que a própria atualização realizada pelo
fluxo gere sucessivas notificações de encerramento, foi
utilizado um campo de controle denominado
``Notificação de fechamento enviada?''. O ramo de
conclusão somente envia o e-mail quando esse campo indica
que a notificação ainda não foi realizada; na sequência, o
campo é atualizado para \textit{Sim}. Esse mecanismo atua
como proteção contra reprocessamento e reduz o risco de
\textit{loop} provocado pelo gatilho de modificação da
lista.

No ramo \textit{Cancelado}, o solicitante recebe uma
notificação específica informando que o chamado foi
cancelado pela Qualidade/SAC. A comunicação apresenta os
dados do chamado e o motivo do cancelamento registrado no
item, mantendo o histórico disponível para consulta.

\subsection{Estados do Chamado e Responsabilidades}

A implementação consolidou cinco estados principais:
\textit{Aberto}, \textit{Em tratativa}, \textit{Em Análise},
\textit{Concluído} e \textit{Cancelado}. A
Tabela~\ref{tab:estados-chamado} apresenta a função de cada
estado.

\begin{table}[H]
\caption{Estados implementados no fluxo de chamados}
\label{tab:estados-chamado}
\centering
\begin{tabular}{p{2.6cm} p{4.2cm} p{5.5cm}}
\hline
\textbf{Status} & \textbf{Responsável predominante} &
\textbf{Significado no processo} \\
\hline
Aberto &
Qualidade/SAC &
Chamado registrado e aguardando classificação e
direcionamento. \\
\hline
Em tratativa &
Área responsável &
Chamado classificado, atribuído e disponível para registro
da causa raiz e da ação corretiva. \\
\hline
Em Análise &
Qualidade/SAC &
Tratativa preenchida pela área responsável e aguardando
validação antes do encerramento. \\
\hline
Concluído &
Qualidade/SAC &
Tratativa validada e chamado encerrado, com notificação aos
envolvidos. \\
\hline
Cancelado &
Qualidade/SAC &
Chamado encerrado sem continuidade da tratativa, com motivo
registrado e comunicação ao solicitante. \\
\hline
\end{tabular}
\raggedright
\par
{\fontsize{10}{12}\selectfont Fonte: Elaborado pelo autor}
\end{table}

\subsection{Rastreabilidade e Notificações}

A rastreabilidade foi implementada por meio da combinação
entre protocolo único, item persistido na Microsoft List,
histórico de modificações e notificações vinculadas às
mudanças de estado. Cada chamado permanece associado ao
mesmo registro durante todo o ciclo, permitindo recuperar
os dados de abertura, classificação, responsável, tratativa
e encerramento.

As notificações foram utilizadas como mecanismo de
transição entre os responsáveis pelo processo. Na abertura,
Qualidade/SAC e solicitante são comunicados. Na atribuição
ou transferência, o responsável pela tratativa é
notificado. Após o registro da causa raiz e da ação
corretiva, a Qualidade/SAC recebe a solicitação para
análise. Na conclusão, solicitante e responsáveis são
informados do encerramento; no cancelamento, o solicitante
recebe a justificativa correspondente.

Esse desenho substitui a comunicação paralela utilizada no
processo anterior por mensagens vinculadas a eventos
registrados na própria base, reduzindo a dependência de
grupos de mensagens para informar movimentações do chamado.

\subsection{Aderência da Implementação aos Requisitos}

A análise entre requisitos e componentes desenvolvidos
demonstrou que a versão atual atende aos elementos centrais
de centralização, protocolo, classificação, atribuição,
transferência, tratativa, validação e encerramento.
Entretanto, alguns requisitos identificados na pesquisa de
campo não puderam ser confirmados como implementados nos
artefatos analisados e devem ser tratados como itens de
evolução ou de validação antes da conclusão definitiva do
projeto.

\begin{table}[H]
\caption{Aderência dos principais requisitos à versão desenvolvida}
\label{tab:aderencia-requisitos}
\centering
\begin{tabular}{p{1.2cm} p{6.6cm} p{4.4cm}}
\hline
\textbf{ID} & \textbf{Requisito} &
\textbf{Situação observada} \\
\hline
RF01--RF14 &
Centralização, abertura, protocolo, notificações,
classificação, atribuição, transferência, tratativa,
análise e encerramento. &
Implementados nos componentes SharePoint, Power Apps,
Microsoft List, FL01 e FL03. \\
\hline
RF15 &
Controle automático de SLA por tipo e prioridade. &
Requisito identificado na pesquisa, porém o cálculo
automático não foi evidenciado nos fluxos exportados
analisados. \\
\hline
RF16 &
Comentários, atualizações e justificativas durante
tratativas prolongadas. &
Necessidade registrada pela Supervisora; não foi
identificado campo específico na estrutura final analisada. \\
\hline
RF17 &
Quantidade de protocolos relacionados. &
Requisito informado pela Analista; não foi identificado no
formulário final de abertura analisado. \\
\hline
RF18 &
Tipologias reutilizáveis. &
Implementado na etapa de classificação da Microsoft List. \\
\hline
RF19 &
Indicadores gerenciais. &
A Central de Chamados possui acesso ao dashboard; a
validação detalhada das métricas deve ser realizada na
etapa de avaliação dos resultados. \\
\hline
\end{tabular}
\raggedright
\par
{\fontsize{10}{12}\selectfont Fonte: Elaborado pelo autor}
\end{table}

Outro ponto de aderência parcial está na regra de
priorização. A pesquisa de campo indicou prioridade Alta
para chamados Massivos e Média para chamados Internos.
Entretanto, na versão exportada do FL01, todo novo chamado
é inicialmente atualizado com prioridade Média, ficando a
classificação definitiva sob responsabilidade da
Qualidade/SAC. Dessa forma, a regra de prioridade por tipo
não está automatizada no momento analisado.

Esses itens não invalidam o fluxo principal desenvolvido,
mas representam diferenças entre os requisitos levantados
e a versão implementada. O registro explícito dessas
diferenças é importante para a avaliação acadêmica da
solução, pois permite distinguir funcionalidades
efetivamente entregues de requisitos planejados ou
reservados para evolução.

\section{Desenvolvimento da Solução}

Esta seção corresponde à etapa de \textit{Automação} do
ciclo de BPM apresentado na Base Teórica: o processo futuro
(TO-BE) especificado na seção de Análise e Redesenho do
Processo é aqui transformado em componentes executáveis do
ecossistema Microsoft 365, preservando, mesmo em ambiente
\textit{low-code}, as etapas de arquitetura, implementação e
validação próprias da Engenharia de Software. O
desenvolvimento foi realizado com base nos requisitos
consolidados na etapa anterior e utilizou componentes do
ecossistema Microsoft 365 já disponível na organização. A
solução resultante foi estruturada em cinco elementos
integrados: uma Central de Chamados no SharePoint, um
aplicativo Canvas no Power Apps para abertura das
solicitações, uma Microsoft List denominada
\textit{Chamados} como repositório central, dois fluxos
principais no Power Automate e uma área de acesso aos
indicadores gerenciais.

A arquitetura implementada difere parcialmente da
arquitetura inicialmente descrita na proposta. Na versão
desenvolvida, o SharePoint atua como portal de acesso,
enquanto o registro do chamado é realizado por um
aplicativo Power Apps incorporado à página. O Power Apps
grava diretamente o novo item na Microsoft List e, a partir
da criação desse item, o fluxo FL01 é disparado para gerar o
protocolo e realizar as notificações. Essa alteração
preservou o SharePoint como ponto único de entrada, porém
separou a camada de interface da camada de armazenamento,
permitindo maior controle sobre a experiência do
solicitante.

\subsection{Arquitetura da Solução Desenvolvida}

A Tabela~\ref{tab:arquitetura-desenvolvida} apresenta os
componentes efetivamente utilizados e suas responsabilidades
na versão desenvolvida.

\begin{table}[H]
\caption{Componentes da arquitetura desenvolvida}
\label{tab:arquitetura-desenvolvida}
\centering
\begin{tabular}{p{3.0cm} p{4.1cm} p{5.1cm}}
\hline
\textbf{Componente} & \textbf{Função principal} &
\textbf{Responsabilidade no fluxo} \\
\hline
SharePoint &
Portal corporativo da solução &
Centralizar o acesso para abertura de chamado, consulta de
registros e acesso ao painel gerencial. \\
\hline
Power Apps &
Interface de abertura &
Coletar os dados do solicitante, validar campos
obrigatórios, permitir anexos e gravar o novo chamado na
Microsoft List. \\
\hline
Microsoft List &
Repositório central &
Armazenar dados de abertura, classificação, responsável,
tratativa, status, datas e histórico de alterações. \\
\hline
Power Automate --- FL01 &
Automação da abertura &
Gerar protocolo, inicializar dados do chamado e notificar
Qualidade/SAC e solicitante. \\
\hline
Power Automate --- FL03 &
Automação da tratativa e encerramento &
Controlar notificações por mudança de status, transferência
de responsável, envio para análise, conclusão e
cancelamento. \\
\hline
Power BI &
Camada gerencial &
Disponibilizar a visualização dos indicadores e da base
histórica a partir do portal da Central de Chamados. \\
\hline
\end{tabular}
\raggedright
\par
{\fontsize{10}{12}\selectfont Fonte: Elaborado pelo autor}
\end{table}

O fluxo lógico da solução pode ser representado pela
sequência: solicitante acessa o SharePoint, realiza a
abertura no Power Apps, o registro é persistido na
Microsoft List, o FL01 gera protocolo e notificações, a
Qualidade/SAC classifica e direciona o chamado, a área
responsável realiza a tratativa e o FL03 controla as
mudanças subsequentes até o encerramento.

\subsection{Central de Chamados no SharePoint}

O SharePoint foi utilizado como camada de apresentação e
ponto central de acesso ao processo. A página desenvolvida
apresenta uma identificação institucional da Central de
Chamados Internos e disponibiliza acessos para abertura de
novo chamado, consulta dos chamados registrados e acesso ao
dashboard gerencial. Essa organização concentra em um único
ambiente as funções que, no processo anterior, estavam
distribuídas entre SAP Service, mensagens e comunicações
paralelas.

A utilização do SharePoint como portal também permite que o
processo seja incorporado ao ambiente corporativo já
utilizado pelos usuários, reduzindo a necessidade de
acesso a uma aplicação externa. O formulário de abertura
foi incorporado à página por meio de um aplicativo Power
Apps, enquanto a Microsoft List permanece responsável pelo
armazenamento e pelas etapas internas de classificação,
tratativa e encerramento.

\subsection{Aplicativo Power Apps para Abertura do Chamado}

A primeira etapa do formulário originalmente prevista
diretamente na Microsoft List foi substituída por um
aplicativo Canvas desenvolvido no Power Apps. A mudança foi
realizada com o objetivo de oferecer uma interface mais
simples ao solicitante e separar os campos de abertura dos
campos internos de classificação e tratamento.

O aplicativo foi integrado diretamente à lista
\textit{Chamados} e configurado para criação de novos
registros. Na versão analisada, o formulário de abertura
contém os campos Empresa, Tipo de chamado, Assunto do
chamado, Problema, Linha, Veículo e Anexos. Empresa, Tipo
de chamado, Assunto do chamado e Problema foram
configurados como obrigatórios. O campo Problema utiliza
entrada de texto em múltiplas linhas, possibilitando uma
descrição detalhada da ocorrência, enquanto o controle de
Anexos permite incluir evidências que apoiem a análise.

Os campos de escolha são carregados a partir das opções
cadastradas na própria Microsoft List, evitando a
manutenção de listas paralelas no aplicativo. O botão
\textit{Criar Chamado} utiliza a validação do formulário
antes do envio, enquanto o botão \textit{Limpar Formulário}
reinicializa os campos para uma nova entrada.

Após o salvamento bem-sucedido, o aplicativo altera seu
estado visual e apresenta ao usuário uma confirmação de que
o chamado foi aberto e está aguardando análise da
Qualidade/SAC. A mesma tela disponibiliza o botão
\textit{Novo Chamado}, que limpa o formulário, retorna ao
modo de criação e permite novo registro. Em caso de falha
no envio, uma mensagem de erro é apresentada ao usuário.

A interface foi ajustada para uso incorporado em página do
SharePoint, utilizando contêineres responsivos e dimensões
baseadas no espaço disponível. Essa adequação buscou
preservar a legibilidade do formulário quando executado
dentro do \textit{web part} do Power Apps.

\subsection{Estrutura da Microsoft List}

A Microsoft List denominada \textit{Chamados} constitui o
repositório central da solução. A estrutura separa os dados
informados pelo solicitante dos dados preenchidos durante o
processamento interno.

Na abertura são armazenados, entre outros, Empresa, Tipo de
chamado, Assunto, Problema, Veículo, Linha, anexos e o
usuário responsável pela criação do item. Nas etapas
posteriores são utilizados os campos Prioridade, Área
Responsável, Tipologia, Responsável pela tratativa, Causa
raiz, Ação corretiva, Data da Tratativa, Status do Chamado,
dados de transferência, Data de Fechamento e controle de
notificação de fechamento.

O registro centralizado permite que todas as etapas
utilizem o mesmo item como referência, evitando a criação
de bases paralelas e mantendo o histórico de alterações
associado ao protocolo. A estrutura também sustenta as
automações do Power Automate, que consultam alterações de
campos específicos para identificar quando uma nova ação
deve ser executada.

\subsection{Fluxo FL01 --- Abertura de Chamados Internos}

O fluxo FL01 é iniciado automaticamente quando um novo item
é criado na Microsoft List. Sua primeira ação calcula o
protocolo utilizando o ano corrente e o identificador
numérico do item. O padrão implementado é
\texttt{C-AAAA-NNNN}. Assim, um item de identificador 14
criado em 2026, por exemplo, recebe o protocolo
\texttt{C-2026-0014}.

Após a geração do protocolo, o próprio item é atualizado.
Na versão exportada do fluxo, o status inicial é definido
como \textit{Aberto} e a prioridade é inicialmente
registrada como \textit{Média}. Em seguida são executadas
duas notificações distintas.

A primeira notificação é destinada à equipe de
Qualidade/SAC e informa que um novo chamado aguarda
classificação. O e-mail apresenta protocolo, solicitante,
empresa, tipo, assunto, linha, veículo, status e descrição
do problema, além de um link direto para o item. O objetivo
é iniciar a etapa de classificação e direcionamento.

A segunda notificação é destinada ao próprio solicitante e
funciona como confirmação formal de abertura. O e-mail
informa o número do protocolo, resume os dados registrados,
indica que a solicitação aguarda análise da Qualidade/SAC e
fornece o link para acompanhamento do chamado. Com isso, a
confirmação de abertura deixa de depender de comunicação
manual.

\subsection{Classificação e Direcionamento pela
Qualidade/SAC}

Após a abertura, a Qualidade/SAC realiza a classificação
diretamente no registro da Microsoft List. Nessa etapa são
definidos os campos Prioridade, Área Responsável, Tipologia
e Responsável pela tratativa. O status do chamado também é
controlado pela Qualidade/SAC.

A decisão de manter esses campos fora do formulário do
solicitante reduz o risco de encaminhamento incorreto e
formaliza a Qualidade/SAC como instância de triagem. Quando
o chamado é validado e direcionado para tratamento, o
status é alterado para \textit{Em tratativa}, iniciando a
execução das regras do FL03 associadas à atribuição do
responsável.

A estrutura da lista também contempla o campo Motivo do
Cancelamento. Quando a Qualidade/SAC decide não prosseguir
com a solicitação e altera o status para
\textit{Cancelado}, o motivo registrado passa a compor a
rastreabilidade do encerramento sem tratativa.

\subsection{Fluxo FL03 --- Tratativa e Encerramento}

O FL03 é disparado quando um item da lista é modificado.
Inicialmente, o fluxo obtém quais propriedades foram
alteradas e recupera o estado atual do item. Em seguida,
uma estrutura de decisão do tipo \textit{Switch} avalia o
status do chamado e direciona o processamento para o ramo
correspondente.

No estado \textit{Em tratativa}, o fluxo verifica se houve
mudança do status para identificar uma nova atribuição.
Quando essa condição é atendida, o responsável pela
tratativa recebe uma notificação por e-mail contendo os
dados necessários e o link para acesso ao chamado.

Ainda no estado \textit{Em tratativa}, o fluxo verifica se
o campo de responsável foi alterado. Quando ocorre uma
transferência, o registro é atualizado com informações de
controle da transferência e o novo responsável é
notificado. Essa regra foi incorporada após a necessidade
identificada no levantamento de requisitos de redistribuir
chamados sem perda de rastreabilidade.

Quando Causa raiz e Ação corretiva são preenchidas e as
respectivas propriedades são identificadas como alteradas,
o FL03 atualiza automaticamente o status para
\textit{Em Análise} e envia uma notificação à
Qualidade/SAC. A área responsável, portanto, registra a
tratativa, mas não realiza o encerramento definitivo.

A análise final é responsabilidade da Qualidade/SAC. Após
avaliar a tratativa, a equipe pode alterar manualmente o
status para \textit{Concluído} ou, quando aplicável,
\textit{Cancelado}. No caso de conclusão, o FL03 notifica o
solicitante e os responsáveis envolvidos e registra a data
de fechamento.

Para impedir que a própria atualização realizada pelo
fluxo gere sucessivas notificações de encerramento, foi
utilizado um campo de controle denominado
``Notificação de fechamento enviada?''. O ramo de
conclusão somente envia o e-mail quando esse campo indica
que a notificação ainda não foi realizada; na sequência, o
campo é atualizado para \textit{Sim}. Esse mecanismo atua
como proteção contra reprocessamento e reduz o risco de
\textit{loop} provocado pelo gatilho de modificação da
lista.

No ramo \textit{Cancelado}, o solicitante recebe uma
notificação específica informando que o chamado foi
cancelado pela Qualidade/SAC. A comunicação apresenta os
dados do chamado e o motivo do cancelamento registrado no
item, mantendo o histórico disponível para consulta.

\subsection{Estados do Chamado e Responsabilidades}

A implementação consolidou cinco estados principais:
\textit{Aberto}, \textit{Em tratativa}, \textit{Em Análise},
\textit{Concluído} e \textit{Cancelado}. A
Tabela~\ref{tab:estados-chamado} apresenta a função de cada
estado.

\begin{table}[H]
\caption{Estados implementados no fluxo de chamados}
\label{tab:estados-chamado}
\centering
\begin{tabular}{p{2.6cm} p{4.2cm} p{5.5cm}}
\hline
\textbf{Status} & \textbf{Responsável predominante} &
\textbf{Significado no processo} \\
\hline
Aberto &
Qualidade/SAC &
Chamado registrado e aguardando classificação e
direcionamento. \\
\hline
Em tratativa &
Área responsável &
Chamado classificado, atribuído e disponível para registro
da causa raiz e da ação corretiva. \\
\hline
Em Análise &
Qualidade/SAC &
Tratativa preenchida pela área responsável e aguardando
validação antes do encerramento. \\
\hline
Concluído &
Qualidade/SAC &
Tratativa validada e chamado encerrado, com notificação aos
envolvidos. \\
\hline
Cancelado &
Qualidade/SAC &
Chamado encerrado sem continuidade da tratativa, com motivo
registrado e comunicação ao solicitante. \\
\hline
\end{tabular}
\raggedright
\par
{\fontsize{10}{12}\selectfont Fonte: Elaborado pelo autor}
\end{table}

\subsection{Rastreabilidade e Notificações}

A rastreabilidade foi implementada por meio da combinação
entre protocolo único, item persistido na Microsoft List,
histórico de modificações e notificações vinculadas às
mudanças de estado. Cada chamado permanece associado ao
mesmo registro durante todo o ciclo, permitindo recuperar
os dados de abertura, classificação, responsável, tratativa
e encerramento.

As notificações foram utilizadas como mecanismo de
transição entre os responsáveis pelo processo. Na abertura,
Qualidade/SAC e solicitante são comunicados. Na atribuição
ou transferência, o responsável pela tratativa é
notificado. Após o registro da causa raiz e da ação
corretiva, a Qualidade/SAC recebe a solicitação para
análise. Na conclusão, solicitante e responsáveis são
informados do encerramento; no cancelamento, o solicitante
recebe a justificativa correspondente.

Esse desenho substitui a comunicação paralela utilizada no
processo anterior por mensagens vinculadas a eventos
registrados na própria base, reduzindo a dependência de
grupos de mensagens para informar movimentações do chamado.

\subsection{Aderência da Implementação aos Requisitos}

A análise entre requisitos e componentes desenvolvidos
demonstrou que a versão atual atende aos elementos centrais
de centralização, protocolo, classificação, atribuição,
transferência, tratativa, validação e encerramento.
Entretanto, alguns requisitos identificados na pesquisa de
campo não puderam ser confirmados como implementados nos
artefatos analisados e devem ser tratados como itens de
evolução ou de validação antes da conclusão definitiva do
projeto.

\begin{table}[H]
\caption{Aderência dos principais requisitos à versão desenvolvida}
\label{tab:aderencia-requisitos}
\centering
\begin{tabular}{p{1.2cm} p{6.6cm} p{4.4cm}}
\hline
\textbf{ID} & \textbf{Requisito} &
\textbf{Situação observada} \\
\hline
RF01--RF14 &
Centralização, abertura, protocolo, notificações,
classificação, atribuição, transferência, tratativa,
análise e encerramento. &
Implementados nos componentes SharePoint, Power Apps,
Microsoft List, FL01 e FL03. \\
\hline
RF15 &
Controle automático de SLA por tipo e prioridade. &
Requisito identificado na pesquisa, porém o cálculo
automático não foi evidenciado nos fluxos exportados
analisados. \\
\hline
RF16 &
Comentários, atualizações e justificativas durante
tratativas prolongadas. &
Necessidade registrada pela Supervisora; não foi
identificado campo específico na estrutura final analisada. \\
\hline
RF17 &
Quantidade de protocolos relacionados. &
Requisito informado pela Analista; não foi identificado no
formulário final de abertura analisado. \\
\hline
RF18 &
Tipologias reutilizáveis. &
Implementado na etapa de classificação da Microsoft List. \\
\hline
RF19 &
Indicadores gerenciais. &
A Central de Chamados possui acesso ao dashboard; a
validação detalhada das métricas deve ser realizada na
etapa de avaliação dos resultados. \\
\hline
\end{tabular}
\raggedright
\par
{\fontsize{10}{12}\selectfont Fonte: Elaborado pelo autor}
\end{table}

Outro ponto de aderência parcial está na regra de
priorização. A pesquisa de campo indicou prioridade Alta
para chamados Massivos e Média para chamados Internos.
Entretanto, na versão exportada do FL01, todo novo chamado
é inicialmente atualizado com prioridade Média, ficando a
classificação definitiva sob responsabilidade da
Qualidade/SAC. Dessa forma, a regra de prioridade por tipo
não está automatizada no momento analisado.

Esses itens não invalidam o fluxo principal desenvolvido,
mas representam diferenças entre os requisitos levantados
e a versão implementada. O registro explícito dessas
diferenças é importante para a avaliação acadêmica da
solução, pois permite distinguir funcionalidades
efetivamente entregues de requisitos planejados ou
reservados para evolução.

\section{Resultados e Discussão}

Apresentar os resultados da pesquisa, podendo utilizar ilustrações como tabelas, figuras, quadros, fluxogramas e imagens.

\textbf{Normas para utilização das ilustrações:}

\begin{itemize}

\item Elas devem possuir títulos (cabeçalhos) localizados na parte superior, antecedidos da palavra que as designa (tabela, figura, esquema, fluxograma, imagem etc.), seguidos do número de ordem de ocorrência no texto, em algarismos arábicos, e travessão para separação do título.

\item As fontes consultadas são obrigatórias, mesmo que sejam de produção do próprio autor, devendo estar localizadas na parte inferior, contendo notas e outras informações necessárias à sua compreensão, quando aplicável.

\item As ilustrações devem ser centralizadas na página.

\item Os títulos e as fontes devem ficar alinhados à borda esquerda da ilustração. Caso ultrapassem sua largura, devem ser centralizados.

\item Deve-se utilizar fonte Arial, tamanho 12 para o título e tamanho 10 para a fonte, ambos sem ponto final.

\item Utilizar o termo \textbf{Figura} para desenhos e \textbf{Imagem} para fotografias.

\item As tabelas devem apresentar dados numéricos (valores, porcentagens, estatísticas etc.), enquanto os quadros devem apresentar dados qualitativos (ideias, descrições, resumos), sem uso de abreviações.

\item Caso a tabela não caiba em uma página, deve-se indicar ``continua'' ao final da primeira página e ``continuação'' no início da página seguinte, repetindo o título e o cabeçalho.

\end{itemize}

\vspace{0.5cm}

\textbf{Exemplo de Figura:}

% ================= FIGURA =================
\begin{figure}[H]
\caption{Figura 1 -- Exemplo de figura}
\centering
\includegraphics[width=0.5\textwidth]{Imagens/exemplo.jpg}

\raggedright
{\fontsize{10}{12}\selectfont Fonte: Elaborado pelo autor}
\end{figure}

\vspace{0.5cm}

\textbf{Exemplo de Tabela:}

% ================= TABELA =================
\begin{table}[H]
\caption{Tabela 1 -- Exemplo de tabela}
\centering
\begin{tabular}{lcc}
\hline
Item & Quantidade & Percentual \\
\hline
Teoria social & 22 & 7,9\% \\
Método & 34 & 12,3\% \\
Questão & 54 & 19,5\% \\
Raciocínio & 124 & 44,8\% \\
\hline
\end{tabular}

\raggedright
{\fontsize{10}{12}\selectfont Fonte: Adaptado de \cite{onu2023}}
\end{table}
\section{Considerações Finais}

Apresentar as considerações finais a partir dos resultados encontrados com a pesquisa. Descrição das principais contribuições do estudo e indicativos para estudos futuros. 

\subsection{AGRADECIMENTOS} %Opcional!
Parágrafo único, espaçamento simples, fonte Arial, tamanho 12, sem negrito, alinhamento do texto justificado, espaçamento simples. 

% ANEXO e/ou APÊNDICE(seção opcional)
% Parágrafo único, espaçamento simples, fonte Arial, tamanho 12, sem negrito, alinhamento do texto justificado, espaçamento simples. 


% ================= REFERÊNCIAS =================
\printbibliography

\end{document}
